\input{preamble}

% OK, start here.
%
\begin{document}

\title{Morphismes \'etales de schémas}

\maketitle

\phantomsection
\label{section-phantom}

\tableofcontents



\section{Introduction}
\label{section-introduction}

\noindent
Dans ce chapitre, nous étudions les morphismes \'etales de schémas. Nous illustrons
certains des concepts les plus importants en travaillant dans le cas noethérien.
Notre objectif principal est de réunir pour le lecteur suffisamment de résultats
d'algèbre commutative pour commencer la lecture d'un traité de cohomologie \'etale. Un
objectif auxiliaire est de fournir assez d'éléments pour que le lecteur puisse
cesser de qualifier l'expression « la topologie \'etale des schémas »
d'exercice de pur formalisme abstrait, s'il se livre à un tel blasphème.

\medskip\noindent
Nous renverrons aux autres
chapitres du projet Stacks pour les résultats classiques de géométrie algébrique
(sur les schémas et l'algèbre commutative). Nous donnerons des démonstrations
détaillées des nouveaux résultats que nous énonçons ici.




\section{Conventions}
\label{section-conventions}

\noindent
Dans ce chapitre, les schémas seront souvent supposés localement noethériens
et les anneaux seront souvent supposés noethériens. Mais, dans tous les énoncés,
nous le répéterons chaque fois que cela sera nécessaire et veillerons à énumérer
toutes les hypothèses ! Voici par ailleurs quelques faits généraux que nous utiliserons
souvent et qu'il est utile de garder à l'esprit :
\begin{enumerate}
\item Un homomorphisme d'anneaux $A \to B$ de type fini, avec $A$ noethérien,
est de présentation finie. Voir Algèbre,
Lemme \ref{algebra-lemma-Noetherian-finite-type-is-finite-presentation}.
\item Un morphisme (localement) de type fini entre schémas localement noethériens
est automatiquement (localement) de présentation finie.
Voir Morphismes,
Lemme \ref{morphisms-lemma-noetherian-finite-type-finite-presentation}.
\item Ajouter ici d'autres faits de ce genre.
\end{enumerate}




\section{Morphismes non ramifiés}
\label{section-unramified-definition}

\noindent
Nous définissons d'abord les « homomorphismes non ramifiés d'anneaux locaux » pour les anneaux
locaux noethériens. Nous ne pouvons employer simplement le terme « non ramifié », car il existe déjà
une notion
d'homomorphisme d'anneaux non ramifié (Algèbre, Section \ref{algebra-section-unramified})
et elle est différente. Après avoir quelque peu étudié cette notion, nous
la globaliserons pour décrire les morphismes non ramifiés de schémas localement noethériens.

\begin{definition}
\label{definition-unramified-rings}
Soient $A$ et $B$ des anneaux locaux noethériens. Un homomorphisme local $A \to B$
est appelé un {\it homomorphisme non ramifié d'anneaux locaux} si
\begin{enumerate}
\item $\mathfrak m_AB = \mathfrak m_B$,
\item $\kappa(\mathfrak m_B)$ est une extension finie séparable de
$\kappa(\mathfrak m_A)$, et
\item $B$ est essentiellement de type fini sur $A$ (cela signifie
que $B$ est le localisé d'une $A$-algèbre de type fini en un idéal premier).
\end{enumerate}
\end{definition}

\noindent
C'est la version locale de la
définition donnée dans Algèbre, Section \ref{algebra-section-unramified}.
Dans cette section, un homomorphisme d'anneaux $R \to S$ est dit non ramifié si et
seulement s'il est de type fini et si $\Omega_{S/R} = 0$.
Nous disons que $R \to S$ est non ramifié en un idéal premier $\mathfrak q \subset S$
s'il existe $g \in S$, $g \not \in \mathfrak q$, tel que
$R \to S_g$ soit un homomorphisme d'anneaux non ramifié. Il est démontré dans
Algèbre, Lemmes \ref{algebra-lemma-unramified-at-prime} et
\ref{algebra-lemma-characterize-unramified}, qu'étant donnés un homomorphisme
d'anneaux $R \to S$ de type fini et un idéal premier $\mathfrak q$ de $S$
au-dessus de $\mathfrak p \subset R$, on a
$$
R \to S\text{ est non ramifié en }\mathfrak q
\Leftrightarrow
\mathfrak pS_{\mathfrak q} = \mathfrak q S_{\mathfrak q}
\text{ et }
\kappa(\mathfrak p) \subset \kappa(\mathfrak q)\text{ est finie séparable}
$$
Nous voyons donc que, pour un homomorphisme local d'anneaux locaux, les propriétés
de notre définition ci-dessus sont étroitement liées au fait
d'être non ramifié. Nous avons en fait démontré le lemme suivant.

\begin{lemma}
\label{lemma-characterize-unramified-Noetherian}
\begin{slogan}
La non-ramification est une condition locale sur les germes.
\end{slogan}
Soit $A \to B$ un homomorphisme de type fini, avec $A$ noethérien.
Soit $\mathfrak q$ un idéal premier de $B$ au-dessus de $\mathfrak p \subset A$.
Alors $A \to B$ est non ramifié en $\mathfrak q$ si et seulement si
$A_{\mathfrak p} \to B_{\mathfrak q}$ est un homomorphisme non ramifié
d'anneaux locaux.
\end{lemma}

\begin{proof}
Voir la discussion précédente.
\end{proof}

\noindent
Nous allons caractériser la propriété d'être non ramifié au moyen
des complétés. Pour un anneau local noethérien $A$,
nous notons $A^\wedge$ le complété de $A$ pour
l'idéal maximal. C'est encore un anneau local noethérien ; voir
Algèbre, Lemme \ref{algebra-lemma-completion-Noetherian-Noetherian}.

\begin{lemma}
\label{lemma-unramified-completions}
Soient $A$ et $B$ des anneaux locaux noethériens.
Soit $A \to B$ un homomorphisme local.
\begin{enumerate}
\item si $A \to B$ est un homomorphisme non ramifié d'anneaux locaux,
alors $B^\wedge$ est un $A^\wedge$-module de type fini,
\item si $A \to B$ est un homomorphisme non ramifié d'anneaux locaux et si
$\kappa(\mathfrak m_A) = \kappa(\mathfrak m_B)$,
alors $A^\wedge \to B^\wedge$ est surjectif,
\item si $A \to B$ est un homomorphisme non ramifié d'anneaux locaux
et si $\kappa(\mathfrak m_A)$
est séparablement clos, alors $A^\wedge \to B^\wedge$ est surjectif,
\item si $A$ et $B$ sont des anneaux de valuation discrète complets, alors
$A \to B$ est un homomorphisme non ramifié d'anneaux locaux
si et seulement si une uniformisante de $A$ s'envoie sur une uniformisante de $B$,
et si l'extension des corps résiduels est finie séparable (et si $B$ est
essentiellement de type fini sur $A$).
\end{enumerate}
\end{lemma}

\begin{proof}
L'assertion (1) est un cas particulier du résultat
Algèbre, Lemme \ref{algebra-lemma-finite-after-completion}.
Pour l'assertion (2), remarquons que l'espace vectoriel sur $\kappa(\mathfrak m_A)$
$B^\wedge/\mathfrak m_{A^\wedge}B^\wedge$
est engendré par $1$. Le lemme de Nakayama
(Algèbre, Lemme \ref{algebra-lemma-NAK}) montre donc que l'application
$A^\wedge \to B^\wedge$ est surjective.
L'assertion (3) est un cas particulier de l'assertion (2).
L'assertion (4) découle immédiatement des définitions.
\end{proof}

\begin{lemma}
\label{lemma-characterize-unramified-completions}
Soient $A$ et $B$ des anneaux locaux noethériens.
Soit $A \to B$ un homomorphisme local tel que $B$ soit
essentiellement de type fini sur $A$.
Les conditions suivantes sont équivalentes :
\begin{enumerate}
\item $A \to B$ est un homomorphisme non ramifié d'anneaux locaux,
\item $A^\wedge \to B^\wedge$ est un homomorphisme non ramifié d'anneaux locaux, et
\item $A^\wedge \to B^\wedge$ est non ramifié.
\end{enumerate}
\end{lemma}

\begin{proof}
L'équivalence de (1) et (2) résulte du fait que
$\mathfrak m_AA^\wedge$ est l'idéal maximal de $A^\wedge$
(et de même pour $B$), ainsi que de la fidèle platitude de $B \to B^\wedge$.
Par exemple, si $A^\wedge \to B^\wedge$ est non ramifié, alors
$\mathfrak m_AB^\wedge = (\mathfrak m_AB)B^\wedge = \mathfrak m_BB^\wedge$
et donc $\mathfrak m_AB = \mathfrak m_B$.

\medskip\noindent
Supposons satisfaites les conditions équivalentes (1) et (2).
Le Lemme \ref{lemma-unramified-completions}
montre que $A^\wedge \to B^\wedge$ est
fini. Ainsi, $A^\wedge \to B^\wedge$ est de présentation finie et, par
Algèbre, Lemme \ref{algebra-lemma-characterize-unramified},
nous concluons que $A^\wedge \to B^\wedge$ est non ramifié en
$\mathfrak m_{B^\wedge}$. Comme $B^\wedge$ est local, nous concluons
que $A^\wedge \to B^\wedge$ est non ramifié.

\medskip\noindent
Supposons (3). Par Algèbre, Lemme \ref{algebra-lemma-unramified-at-prime},
nous concluons que $A^\wedge \to B^\wedge$ est un homomorphisme non ramifié
d'anneaux locaux, c'est-à-dire que (2) est satisfaite.
\end{proof}

\begin{definition}
\label{definition-unramified-schemes}
(Voir Morphismes, Définition \ref{morphisms-definition-unramified},
pour la définition dans le cas général.)
Soit $Y$ un schéma localement noethérien.
Soit $f : X \to Y$ un morphisme localement de type fini.
Soit $x \in X$.
\begin{enumerate}
\item Nous disons que $f$ est {\it non ramifié en $x$} si
$\mathcal{O}_{Y, f(x)} \to \mathcal{O}_{X, x}$
est un homomorphisme non ramifié d'anneaux locaux.
\item Le morphisme $f : X \to Y$ est dit {\it non ramifié}
s'il est non ramifié en tout point de $X$.
\end{enumerate}
\end{definition}

\noindent
Montrons que cette définition coïncide avec celle du
chapitre sur les morphismes de schémas. Cela garantit en particulier que
l'ensemble des points où un morphisme est non ramifié est ouvert.

\begin{lemma}
\label{lemma-unramified-definition}
Soit $Y$ un schéma localement noethérien.
Soit $f : X \to Y$ un morphisme localement de type fini.
Soit $x \in X$. Le morphisme $f$ est non ramifié en $x$ au
sens de la Définition \ref{definition-unramified-schemes}
si et seulement s'il est non ramifié au
sens de Morphismes, Définition \ref{morphisms-definition-unramified}.
\end{lemma}

\begin{proof}
Cela résulte du Lemme \ref{lemma-characterize-unramified-Noetherian}
et des définitions.
\end{proof}

\noindent
Voici quelques résultats sur les morphismes non ramifiés.
Sous la forme donnée dans cette liste, ils ne s'appliquent qu'aux
morphismes localement de type fini entre schémas localement noethériens.
Dans chaque cas, nous renvoyons au résultat général
démontré plus tôt dans le projet, même si, dans certains cas, on peut
démontrer le résultat plus facilement dans le cas noethérien.
En voici la liste :
\begin{enumerate}
\item La non-ramification est locale sur la source et sur le but pour la
topologie de Zariski.
\item Les morphismes non ramifiés sont stables par changement de base et par composition.
Voir Morphismes, Lemmes \ref{morphisms-lemma-base-change-unramified}
et \ref{morphisms-lemma-composition-unramified}.
\item Les morphismes non ramifiés de schémas sont localement quasi-finis,
et les morphismes non ramifiés quasi-compacts sont quasi-finis.
Voir Morphismes, Lemme \ref{morphisms-lemma-unramified-quasi-finite}.
\item Les morphismes non ramifiés sont de dimension relative $0$. Voir
Morphismes, Définition \ref{morphisms-definition-relative-dimension-d}
et
Morphismes, Lemme \ref{morphisms-lemma-locally-quasi-finite-rel-dimension-0}.
\item Un morphisme est non ramifié si et seulement si toutes ses fibres sont non ramifiées.
Autrement dit, la non-ramification se vérifie sur les fibres schématiques. Voir
Morphismes, Lemme \ref{morphisms-lemma-unramified-etale-fibres}.
\item Soient $X$ et $Y$ non ramifiés sur un schéma de base $S$.
Tout $S$-morphisme de $X$ vers $Y$ est non ramifié.
Voir Morphismes, Lemme \ref{morphisms-lemma-unramified-permanence}.
\end{enumerate}

\section{Trois autres caractérisations des morphismes non ramifiés}
\label{section-three-other}

\noindent
Le théorème suivant donne trois notions équivalentes de
non-ramification en un point. Voir
Morphismes, Lemme \ref{morphisms-lemma-unramified-at-point},
pour une partie de l'énoncé concernant les schémas généraux.

\begin{theorem}
\label{theorem-unramified-equivalence}
Soit $Y$ un schéma localement noethérien.
Soit $f : X \to Y$ un morphisme de schémas localement de type fini.
Soit $x$ un point de $X$. Les conditions suivantes sont équivalentes :
\begin{enumerate}
\item $f$ est non ramifié en $x$,
\item le germe $\Omega_{X/Y, x}$ du module des différentielles relatives
en $x$ est nul,
\item il existe des voisinages ouverts $U$ de $x$ et $V$ de $f(x)$, ainsi qu'un
diagramme commutatif
$$
\xymatrix{
U \ar[rr]_i \ar[rd] & & \mathbf{A}^n_V \ar[ld] \\
& V
}
$$
où $i$ est une immersion fermée définie par un
faisceau quasi-cohérent d'idéaux $\mathcal{I}$ tel que les différentielles
$\text{d}g$, pour $g \in \mathcal{I}_{i(x)}$, engendrent
$\Omega_{\mathbf{A}^n_V/V, i(x)}$, et
\item la diagonale $\Delta_{X/Y} : X \to X \times_Y X$
est un isomorphisme local en $x$.
\end{enumerate}
\end{theorem}

\begin{proof}
L'équivalence de (1) et (2) est démontrée dans
Morphismes, Lemme \ref{morphisms-lemma-unramified-at-point}.

\medskip\noindent
Si $f$ est non ramifié en $x$, alors $f$ est non ramifié sur un voisinage
ouvert de $x$ ; cela ne résulte pas immédiatement
de la Définition \ref{definition-unramified-schemes} de ce chapitre,
mais cela découle de
Morphismes, Définition \ref{morphisms-definition-unramified}, dont nous
avons démontré l'équivalence au
Lemme \ref{lemma-unramified-definition}.
Choisissons des ouverts affines $V \subset Y$ et $U \subset X$
tels que $f(U) \subset V$ et $x \in U$, et tels que $f$ soit
non ramifié sur $U$, c'est-à-dire que $f|_U : U \to V$ soit non ramifié.
Par Morphismes, Lemme \ref{morphisms-lemma-diagonal-unramified-morphism},
le morphisme $U \to U \times_V U$
est une immersion ouverte. Cela démontre que (1) implique (4).

\medskip\noindent
Si $\Delta_{X/Y}$ est un isomorphisme local en $x$, alors
$\Omega_{X/Y, x} = 0$ par
Morphismes, Lemme \ref{morphisms-lemma-differentials-diagonal}.
Nous voyons donc que (4) implique (2).
À ce stade, nous savons que (1), (2) et (4) sont toutes équivalentes.

\medskip\noindent
Supposons (3). L'hypothèse sur le diagramme, combinée au résultat
Morphismes, Lemme \ref{morphisms-lemma-differentials-relative-immersion},
montre que $\Omega_{U/V, x} = 0$. Comme $\Omega_{U/V, x} = \Omega_{X/Y, x}$,
nous concluons que (2) est satisfaite.

\medskip\noindent
Supposons enfin que (2) soit satisfaite. Pour démontrer (3), nous pouvons localiser sur
$X$ et $Y$ et supposer que $X$ et $Y$ sont affines.
Écrivons $X = \Spec(B)$ et $Y = \Spec(A)$.
Le point $x \in X$ correspond à un idéal premier $\mathfrak q \subset B$.
Notre hypothèse est que $\Omega_{B/A, \mathfrak q} = 0$
(voir Morphismes, Lemme \ref{morphisms-lemma-differentials-affine}, pour le
lien entre les différentielles sur les schémas et les modules
de différentielles en algèbre commutative).
Comme $Y$ est localement noethérien et $f$ localement de type fini,
nous voyons que $A$ est noethérien et que
$B \cong A[x_1, \ldots, x_n]/(f_1, \ldots, f_m)$ ; voir
Propriétés, Lemme \ref{properties-lemma-locally-Noetherian}, et
Morphismes, Lemme \ref{morphisms-lemma-locally-finite-type-characterize}.
En particulier, $\Omega_{B/A}$ est un $B$-module de type fini. Nous
pouvons donc trouver un seul $g \in B$, $g \not \in \mathfrak q$, tel que
le localisé principal $(\Omega_{B/A})_g$ soit nul. Ainsi, après
avoir remplacé $B$ par $B_g$, nous voyons que $\Omega_{B/A} = 0$ (la formation
des modules de différentielles commute à la localisation ; voir
Algèbre, Lemme \ref{algebra-lemma-differentials-localize}). Cela signifie que
les $\text{d}(f_j)$ engendrent le noyau de l'application canonique
$\Omega_{A[x_1, \ldots, x_n]/A} \otimes_A B \to \Omega_{B/A}$.
Ainsi, la surjection $A[x_1, \ldots, x_n] \to B$ de $A$-algèbres fournit le
diagramme commutatif de (3), et le théorème est démontré.
\end{proof}

\noindent
Comment utiliser ce théorème ? Voici quelques remarques :
\begin{enumerate}
\item Supposons que
$f : X \to Y$ et $g : Y \to Z$ soient deux morphismes localement de type
fini entre schémas localement noethériens. Il existe une suite
suite exacte courte canonique
$$
f^*(\Omega_{Y/Z}) \to \Omega_{X/Z} \to \Omega_{X/Y} \to 0
$$
voir Morphismes, Lemme \ref{morphisms-lemma-triangle-differentials}.
Le théorème implique donc que, si $g \circ f$ est non ramifié,
alors $f$ l'est aussi. C'est
Morphismes, Lemme \ref{morphisms-lemma-unramified-permanence}.
\item Puisque $\Omega_{X/Y}$ est isomorphe au faisceau conormal
du morphisme diagonal
(Morphismes, Lemme \ref{morphisms-lemma-differentials-diagonal}),
nous voyons que, si $X \to Y$ est un monomorphisme de
schémas localement noethériens et est localement de type fini,
alors $X \to Y$ est non ramifié.
En particulier, les immersions ouvertes et fermées de schémas localement noethériens
sont non ramifiées. Voir
Morphismes, Lemmes
\ref{morphisms-lemma-open-immersion-unramified} et
\ref{morphisms-lemma-closed-immersion-unramified}.
\item Le théorème implique aussi que l'ensemble des points
où un morphisme $f : X \to Y$ (localement de type fini entre schémas localement
noethériens) n'est pas non ramifié est
le support du faisceau cohérent $\Omega_{X/Y}$.
Cela permet de donner une définition schématique du
« lieu de ramification ».
\end{enumerate}

\section{La caractérisation fonctorielle des morphismes non ramifiés}
\label{section-functorial-unramified}

\noindent
En géométrie algébrique élémentaire, on apprend que certaines classes de morphismes peuvent être
caractérisées fonctoriellement et que de telles descriptions sont fort utiles.
Les morphismes non ramifiés admettent eux aussi une telle caractérisation.

\begin{theorem}
\label{theorem-formally-unramified}
Soit $f : X \to S$ un morphisme de schémas.
Supposons que $S$ soit un schéma localement noethérien et que $f$ soit localement de type fini.
Les conditions suivantes sont alors équivalentes :
\begin{enumerate}
\item $f$ est non ramifié,
\item le morphisme $f$ est formellement non ramifié :
pour tout $S$-schéma affine $T$ et tout sous-schéma $T_0$ de $T$
défini par un idéal de carré nul,
l'application naturelle
$$
\Hom_S(T, X) \longrightarrow \Hom_S(T_0, X)
$$
est injective.
\end{enumerate}
\end{theorem}

\begin{proof}
Voir Compléments sur les morphismes,
Lemme \ref{more-morphisms-lemma-unramified-formally-unramified},
pour un énoncé et une démonstration plus généraux.
Voici une esquisse de la démonstration dans le cas présent.

\medskip\noindent
On vérifie d'abord que les deux propriétés sont locales sur la source et sur le but.
Nous pouvons donc supposer que $S$ et $X$ sont affines.
Écrivons $X = \Spec(B)$ et $S = \Spec(R)$.
Écrivons $T = \Spec(C)$. Soit $J$ l'idéal de carré nul de $C$
tel que $T_0 = \Spec(C/J)$. Supposons donné le diagramme
$$
\xymatrix{
& B \ar[d]^\phi \ar[rd]^{\bar{\phi}}
& \\
R \ar[r] \ar[ur] & C \ar[r]
& C/J
}
$$
On vérifie ensuite que l'association $\phi' \mapsto \phi' - \phi$
définit une bijection entre l'ensemble des relèvements de $\bar{\phi}$ et le module
$\text{Der}_R(B, J)$. Ainsi, nous obtenons l'implication (1) $\Rightarrow$ (2)
par la description des morphismes non ramifiés dont le module
de différentielles est nul ; voir le Théorème \ref{theorem-unramified-equivalence}.

\medskip\noindent
Pour obtenir l'implication réciproque, considérons la surjection
$q : C = (B \otimes_R B)/I^2 \to B = C/J$ définie par l'idéal de carré nul
$J = I/I^2$, où $I$ est le noyau de l'application de multiplication
$B \otimes_R B \to B$. Nous disposons déjà d'un relèvement $B \to C$ défini, par exemple,
par $b \mapsto b \otimes 1$. Ainsi, par le même raisonnement que ci-dessus, nous obtenons une
correspondance bijective entre les relèvements de $\text{id} : B \to C/J$ et
$\text{Der}_R(B, J)$. L'hypothèse implique donc que ce dernier module est
nul. Or nous savons que $J \cong \Omega_{B/R}$. Ainsi, $B/R$ est non ramifié.
\end{proof}



\section{Propriétés topologiques des morphismes non ramifiés}
\label{section-topological-unramified}

\noindent
Le premier résultat topologique qui nous sera utile affirme
que les morphismes non ramifiés et séparés ont de « bonnes » sections.
Les résultats de cette section ne nécessitent aucune hypothèse noethérienne.

\begin{proposition}
\label{proposition-properties-sections}
Sections des morphismes non ramifiés.
\begin{enumerate}
\item Toute section d'un morphisme non ramifié est une immersion ouverte.
\item Toute section d'un morphisme séparé est une immersion fermée.
\item Toute section d'un morphisme non ramifié et séparé est ouverte et fermée.
\end{enumerate}
\end{proposition}

\begin{proof}
Fixons un schéma de base $S$.
Si $f : X' \to X$ est un $S$-morphisme quelconque, alors le graphe
$\Gamma_f : X' \to X' \times_S X$
s'obtient par changement de base de la diagonale
$\Delta_{X/S} : X \to X \times_S X$ le long de la projection
$X' \times_S X \to X \times_S X$.
Si $g : X \to S$ est séparé (resp. non ramifié),
alors la diagonale est une immersion fermée (resp. une immersion ouverte)
par Schémas, Définition \ref{schemes-definition-separated}
(resp.\ Morphismes, Lemme \ref{morphisms-lemma-diagonal-unramified-morphism}).
Il en va donc de même du graphe par changement de base (voir
Schémas, Lemme \ref{schemes-lemma-base-change-immersion}).
Dans le cas particulier $X' = S$, nous obtenons (1), resp.\ (2).
L'assertion (3) résulte de la combinaison de (1) et (2).
\end{proof}

\noindent
Nous pouvons maintenant décrire explicitement les sections des morphismes non ramifiés.

\begin{theorem}
\label{theorem-sections-unramified-maps}
Soit $Y$ un schéma connexe.
Soit $f : X \to Y$ un morphisme non ramifié et séparé.
Toute section de $f$ est un isomorphisme sur une composante connexe.
Il existe une correspondance bijective
$$
\text{sections de }f
\leftrightarrow
\left\{
\begin{matrix}
\text{composantes connexes }X'\text{ de }X\text{ telles que}\\
\text{l'application induite }X' \to Y\text{ soit un isomorphisme}
\end{matrix}
\right\}
$$
En particulier, pour $x \in X$, il existe au plus une
section passant par $x$.
\end{theorem}

\begin{proof}
Cela découle directement de la Proposition \ref{proposition-properties-sections}, assertion (3).
\end{proof}

\noindent
Le théorème précédent donne une idée de la « rigidité » des morphismes non ramifiés.
La proposition suivante en donne une autre manifestation ;
outre son intérêt propre, elle est également utile dans la
théorie du groupe fondamental algébrique (voir \cite[Expos\'e V]{SGA1}).
Voir aussi le résultat plus général
Morphismes, Lemme \ref{morphisms-lemma-value-at-one-point}.

\begin{proposition}
\label{proposition-equality}
Soit $S$ un schéma.
Soit $\pi : X \to S$ un morphisme non ramifié et séparé.
Soit $Y$ un $S$-schéma et soit $y \in Y$ un point.
Soient $f, g : Y \to X$ deux $S$-morphismes. Supposons que
\begin{enumerate}
\item $Y$ soit connexe,
\item $x = f(y) = g(y)$, et
\item les applications induites $f^\sharp, g^\sharp : \kappa(x) \to \kappa(y)$
sur les corps résiduels soient égales.
\end{enumerate}
Alors $f = g$.
\end{proposition}

\begin{proof}
Les applications $f, g : Y \to X$ définissent des applications $f', g' : Y \to X_Y = Y \times_S X$
qui sont des sections du morphisme structural $X_Y \to Y$.
Remarquons que $f = g$ si et seulement si $f' = g'$.
Le morphisme structural $X_Y \to Y$ est le changement de base de $\pi$ et est donc
lui aussi non ramifié et séparé (voir
Morphismes, Lemmes \ref{morphisms-lemma-base-change-unramified} et
Schémas, Lemme \ref{schemes-lemma-separated-permanence}).
Ainsi, d'après le Théorème \ref{theorem-sections-unramified-maps},
il suffit de démontrer que $f'$ et $g'$ passent par le même
point de $X_Y$. C'est précisément ce que les hypothèses (2) et (3)
garantissent, à savoir $f'(y) = g'(y) \in X_Y$.
\end{proof}

\begin{lemma}
\label{lemma-finitely-many-maps-to-unramified}
Soit $S$ un schéma noethérien. Soit $X \to S$ un morphisme non ramifié quasi-compact.
Soit $Y \to S$ un morphisme avec $Y$ noethérien. Alors
$\Mor_S(Y, X)$ est un ensemble fini.
\end{lemma}

\begin{proof}
Supposons d'abord que $X \to S$ soit séparé (ce qui est souvent le cas en pratique).
Puisque $Y$ est noethérien, il possède un nombre fini de composantes connexes. Nous
pouvons donc supposer $Y$ connexe. Choisissons un point $y \in Y$, d'image $s \in S$.
Comme $X \to S$ est non ramifié et quasi-compact,
la fibre $X_s$ est finie, disons $X_s = \{x_1, \ldots, x_n\}$,
et $\kappa(x_i)/\kappa(s)$ est une extension finie de corps.
Voir Morphismes, Lemmes \ref{morphisms-lemma-unramified-quasi-finite},
\ref{morphisms-lemma-residue-field-quasi-finite} et
\ref{morphisms-lemma-quasi-finite}.
Pour chaque $i$, il n'existe qu'un nombre fini d'homomorphismes de $\kappa(s)$-algèbres
$\kappa(x_i) \to \kappa(y)$ (par la théorie élémentaire des corps).
Ainsi, $\Mor_S(Y, X)$ est fini par
la Proposition \ref{proposition-equality}.

\medskip\noindent
Cas général. Il existe un ouvert non vide $U \subset S$ tel
que $X_U \to U$ soit fini (donc, en particulier, séparé) ; voir
Morphismes, Lemme \ref{morphisms-lemma-generically-finite}
(ce lemme s'applique puisque nous avons déjà vu ci-dessus qu'un morphisme non ramifié
quasi-compact est quasi-fini et puisque $X \to S$ est quasi-séparé par
Morphismes, Lemme \ref{morphisms-lemma-finite-type-Noetherian-quasi-separated}).
Soit $Z \subset S$ le sous-schéma fermé réduit dont le support est
le complément de $U$. Par récurrence noethérienne, nous voyons que
$\Mor_Z(Y_Z, X_Z)$ est fini (détails omis).
D'après le résultat du premier paragraphe, l'ensemble
$\Mor_U(Y_U, X_U)$ est fini. Il suffit donc de montrer que
$$
\Mor_S(Y, X) \longrightarrow \Mor_Z(Y_Z, X_Z) \times \Mor_U(Y_U, X_U)
$$
est injective. Cela résulte du fait que le lieu où
deux morphismes $a, b : Y \to X$ coïncident est ouvert dans $Y$, puisque
la diagonale $\Delta : X \to X \times_S X$ est ouverte ; voir
Morphismes, Lemme \ref{morphisms-lemma-diagonal-unramified-morphism}.
\end{proof}







\section{Morphismes universellement injectifs et non ramifiés}
\label{section-universally-injective-unramified}

\noindent
Rappelons qu'un morphisme de schémas $f : X \to Y$ est universellement
injectif si tout changement de base de $f$ est injectif (sur les espaces
topologiques sous-jacents) ; voir
Morphismes, Définition \ref{morphisms-definition-universally-injective}.
Les morphismes universellement injectifs et non ramifiés peuvent être
caractérisés comme suit.

\begin{lemma}
\label{lemma-universally-injective-unramified}
Soit $f : X \to S$ un morphisme de schémas.
Les conditions suivantes sont équivalentes :
\begin{enumerate}
\item $f$ est non ramifié et est un monomorphisme,
\item $f$ est non ramifié et universellement injectif,
\item $f$ est localement de type fini et est un monomorphisme,
\item $f$ est universellement injectif, localement de type fini et
formellement non ramifié,
\item $f$ est localement de type fini et $X_s$ est vide
ou bien $X_s \to s$ est un isomorphisme, pour tout $s \in S$.
\end{enumerate}
\end{lemma}

\begin{proof}
Nous avons vu dans
Compléments sur les morphismes, Lemme
\ref{more-morphisms-lemma-unramified-formally-unramified},
qu'être formellement non ramifié et localement de type fini équivaut
à être non ramifié. Ainsi, (4) équivaut à (2).
Un monomorphisme est certainement universellement injectif et
formellement non ramifié ; (3) implique donc (4).
Il est clair que (1) implique (3). Enfin, si (2) est satisfaite, alors
$\Delta : X \to X \times_S X$ est à la fois une immersion ouverte
(Morphismes, Lemme \ref{morphisms-lemma-diagonal-unramified-morphism})
et surjective
(Morphismes, Lemme \ref{morphisms-lemma-universally-injective}) ;
c'est donc un isomorphisme, autrement dit $f$ est un monomorphisme. Ainsi,
(2) implique (1).

\medskip\noindent
La condition (3) implique (5), car les monomorphismes sont préservés par
changement de base
(Schémas, Lemme \ref{schemes-lemma-base-change-monomorphism})
et en vertu de la description des monomorphismes à valeurs dans les spectres de corps
donnée dans
Schémas, Lemme \ref{schemes-lemma-mono-towards-spec-field}.
La condition (5) implique (4) par
Morphismes, Lemmes \ref{morphisms-lemma-universally-injective} et
\ref{morphisms-lemma-unramified-etale-fibres}.
\end{proof}

\noindent
On en déduit la caractérisation utile suivante des immersions fermées.

\begin{lemma}
\label{lemma-characterize-closed-immersion}
Soit $f : X \to S$ un morphisme de schémas.
Les conditions suivantes sont équivalentes :
\begin{enumerate}
\item $f$ est une immersion fermée,
\item $f$ est un monomorphisme propre,
\item $f$ est propre, non ramifié et universellement injectif,
\item $f$ est universellement fermé, non ramifié et est un monomorphisme,
\item $f$ est universellement fermé, non ramifié et universellement injectif,
\item $f$ est universellement fermé, localement de type fini et est un monomorphisme,
\item $f$ est universellement fermé, universellement injectif, localement de
type fini et formellement non ramifié.
\end{enumerate}
\end{lemma}

\begin{proof}
L'équivalence des conditions (4) à (7) résulte immédiatement du
Lemme \ref{lemma-universally-injective-unramified}.

\medskip\noindent
Soit $f : X \to S$ un morphisme qui satisfait à (6). Alors $f$ est séparé ; voir
Schémas, Lemme \ref{schemes-lemma-monomorphism-separated},
et ses fibres sont finies. Ainsi,
Compléments sur les morphismes, Lemme \ref{more-morphisms-lemma-characterize-finite},
montre que $f$ est fini. Puis
Morphismes, Lemme \ref{morphisms-lemma-finite-monomorphism-closed},
implique que $f$ est une immersion fermée, c'est-à-dire que (1) est satisfaite.

\medskip\noindent
Remarquons que (1) $\Rightarrow$ (2), car une immersion fermée est
propre et est un monomorphisme
(Morphismes, Lemme \ref{morphisms-lemma-closed-immersion-proper},
et
Schémas, Lemme \ref{schemes-lemma-immersions-monomorphisms}).
Par le
Lemme \ref{lemma-universally-injective-unramified},
nous voyons que (2) implique (3). Il est clair que (3) implique (5).
\end{proof}

\noindent
Voici un autre résultat du même genre.

\begin{lemma}
\label{lemma-finite-unramified-one-point}
Soit $\pi : X \to S$ un morphisme de schémas. Soit $s \in S$.
Supposons que
\begin{enumerate}
\item $\pi$ soit fini,
\item $\pi$ soit non ramifié,
\item $\pi^{-1}(\{s\}) = \{x\}$, et
\item $\kappa(s) \subset \kappa(x)$ soit purement
inséparable\footnote{Compte tenu de la condition (2),
cela équivaut à $\kappa(s) = \kappa(x)$.}.
\end{enumerate}
Alors il existe un voisinage ouvert $U$ de $s$ tel que
$\pi|_{\pi^{-1}(U)} : \pi^{-1}(U) \to U$ soit une immersion fermée.
\end{lemma}

\begin{proof}
La question est locale sur $S$. Nous pouvons donc supposer que $S = \Spec(A)$.
Par définition d'un morphisme fini, cela implique que $X = \Spec(B)$.
Notons que l'homomorphisme d'anneaux $\varphi : A \to B$ définissant $\pi$
est un homomorphisme d'anneaux fini non ramifié.
Soit $\mathfrak p \subset A$ l'idéal premier correspondant à $s$.
Soit $\mathfrak q \subset B$ l'idéal premier correspondant à $x$.
Les conditions (2), (3) et (4) impliquent que
$B_{\mathfrak q}/\mathfrak pB_{\mathfrak q} = \kappa(\mathfrak p)$.
D'après Algèbre, Lemme \ref{algebra-lemma-unique-prime-over-localize-below},
nous avons $B_{\mathfrak q} = B_{\mathfrak p}$
(notons qu'un homomorphisme d'anneaux fini vérifie la propriété de montée, voir
Algèbre, Section \ref{algebra-section-going-up}.)
Nous voyons donc que
$B_{\mathfrak p}/\mathfrak pB_{\mathfrak p} = \kappa(\mathfrak p)$.
Comme $B$ est un $A$-module fini, le lemme de Nakayama (voir
Algèbre, Lemme \ref{algebra-lemma-NAK})
montre que $B_{\mathfrak p} = \varphi(A_{\mathfrak p})$. Ainsi (en utilisant de
nouveau la finitude de $B$ comme $A$-module), il existe
$f \in A$, $f \not \in \mathfrak p$ tel que $B_f = \varphi(A_f)$,
comme voulu.
\end{proof}

\noindent
Les résultats topologiques présentés ci-dessus serviront à donner une
caractérisation fonctorielle des morphismes \'etales analogue au Théorème
\ref{theorem-formally-unramified}.




\section{Exemples de morphismes non ramifiés}
\label{section-examples}

\noindent
Voici quelques exemples.

\begin{example}
\label{example-etale-field-extensions}
Soit $k$ un corps.
Les morphismes quasi-compacts non ramifiés $X \to \Spec(k)$ sont affines.
En effet, $X$ est de dimension $0$ et noethérien,
donc c'est un ensemble discret fini, et chaque point fournit un ouvert affine ;
ainsi $X$ est une réunion disjointe finie de schémas affines, donc est affine.
Le lemme de normalisation de Noether entraîne que $X$ est le spectre d'une
$k$-algèbre finie $A$.
Cette algèbre est un produit d'extensions finies séparables de $k$.
Ainsi, un morphisme quasi-compact non ramifié vers $\Spec(k)$
correspond à un nombre fini d'extensions finies séparables de $k$.
En particulier, un morphisme non ramifié dont la source est connexe et la cible
réduite à un point correspond nécessairement à une extension finie séparable de corps.
Comme nous le verrons plus loin, $X \to \Spec(k)$ est \'etale si et
seulement s'il est non ramifié. Ainsi, dans ce cas au moins, nous obtenons une
description très simple de la topologie \'etale d'un schéma. Bien entendu, la
cohomologie de cette topologie est une autre affaire.
\end{example}

\begin{example}
\label{example-standard-etale}
La propriété (3) du
Théorème \ref{theorem-unramified-equivalence}
nous fournit une source canonique d'exemples de morphismes non ramifiés.
Fixons un anneau $R$ et un entier $n$. Soit $I = (g_1, \ldots, g_m)$ un
idéal de $R[x_1, \ldots, x_n]$. Soit $\mathfrak q \subset R[x_1, \ldots, x_n]$
un idéal premier. Supposons $I \subset \mathfrak q$ et que la matrice
$$
\left(\frac{\partial g_i}{\partial x_j}\right) \bmod \mathfrak q
\quad\in\quad
\text{Mat}(n \times m, \kappa(\mathfrak q))
$$
soit de rang $n$. Alors le morphisme
$f : Z = \Spec(R[x_1, \ldots, x_n]/I) \to \Spec(R)$
est non ramifié au point $x \in Z \subset \mathbf{A}^n_R$ correspondant
à $\mathfrak q$. Il est clair que nous devons avoir $m \geq n$.
Dans le cas extrême $m = n$, c'est-à-dire lorsque la différentielle de l'application
$\mathbf{A}^n_R \to \mathbf{A}^n_R$ définie par les $g_i$
est un isomorphisme des espaces tangents, alors $f$ est aussi plat
$x$ et est, par conséquent, une application \'etale (voir Algèbre,
Définition \ref{algebra-definition-standard-smooth},
Lemme \ref{algebra-lemma-standard-smooth} et
Exemple \ref{algebra-example-make-standard-smooth}).
\end{example}

\begin{example}
\label{example-number-theory-etale}
Fixons une extension de corps de nombres $L/K$, dont les anneaux d'entiers sont
$\mathcal{O}_L$ et $\mathcal{O}_K$. L'injection $K \to L$ définit un
morphisme $f : \Spec(\mathcal{O}_L) \to \Spec(\mathcal{O}_K)$.
Comme expliqué ci-dessus, les points où $f$ est non ramifié en notre sens
sont exactement ceux où $f$ est non ramifié au sens
usuel. Au sens usuel, le lieu de ramification dans
$\Spec(\mathcal{O}_L)$ est le lieu des zéros de la
différente, qui est un idéal de $\mathcal{O}_L$. En fait, la différente
n'est autre que l'annulateur du module
$\Omega_{\mathcal{O}_L/\mathcal{O}_K}$. De même, le
discriminant est un idéal de $\mathcal{O}_K$, à savoir la
norme de la différente.
Le lieu des zéros du discriminant est précisément l'ensemble
des points de $K$ qui se ramifient dans $L$.
Ainsi, en désignant par $X$ le complémentaire du sous-ensemble fermé
défini par la différente dans $\Spec(\mathcal{O}_L)$,
nous obtenons un morphisme $X \to \Spec(\mathcal{O}_K)$ qui est non ramifié.
De plus, ce morphisme est aussi plat, puisque tout homomorphisme local
d'anneaux de valuation discrète est plat, et ce morphisme est donc
en fait \'etale. Si $L/K$ est une extension galoisienne finie, en désignant par
$Y$ le complémentaire du sous-ensemble fermé défini par le discriminant dans
$\Spec(\mathcal{O}_K)$, nous voyons que nous obtenons même un
morphisme fini \'etale $X \to Y$.
C'est donc un exemple de revêtement \'etale fini.
\end{example}





\section{Morphismes plats}
\label{section-flat-morphisms}

\noindent
Cette section a simplement pour objet de résumer les propriétés de la platitude qui
nous seront utiles. Nous nous contenterons donc d'énoncer précisément les théorèmes
et de donner des références pour leurs démonstrations.

\medskip\noindent
Après avoir brièvement rappelé les faits nécessaires concernant les modules plats sur les anneaux
noethériens, nous énonçons un théorème de Grothendieck qui donne des conditions suffisantes
pour que les « sections hyperplanes » de certains modules soient plates.

\begin{definition}
\label{definition-flat-rings}
Platitude des modules et des anneaux.
\begin{enumerate}
\item Un module $N$ sur un anneau $A$ est dit {\it plat}
si le foncteur $M \mapsto M \otimes_A N$ est exact.
\item Si ce foncteur est également fidèle, nous disons que
$N$ est {\it fidèlement plat} sur $A$.
\item Un morphisme d'anneaux $f : A \to B$ est dit
{\it plat (resp. fidèlement plat)}
si le foncteur $M \mapsto M \otimes_A B$ est exact
(resp. fidèle et exact).
\end{enumerate}
\end{definition}

\noindent
Voici une liste de faits, accompagnés de références au chapitre d'algèbre.
\begin{enumerate}
\item Les modules libres et projectifs sont plats. Cela est clair pour les modules libres
et s'en déduit pour les modules projectifs, car ceux-ci sont facteurs directs de modules
libres et $\otimes$ commute aux sommes directes.
\item La platitude est une propriété locale, c'est-à-dire que $M$ est plat sur $A$
si et seulement si $M_{\mathfrak p}$ est plat sur $A_{\mathfrak p}$ pour tout
$\mathfrak p \in \Spec(A)$.
Voir Algèbre, Lemme \ref{algebra-lemma-flat-localization}.
\item Si $M$ est un $A$-module plat et $A \to B$ un homomorphisme d'anneaux,
alors $M \otimes_A B$ est un $B$-module plat. Voir
Algèbre, Lemme \ref{algebra-lemma-flat-base-change}.
\item Les modules finis et plats sur les anneaux locaux sont libres.
Voir Algèbre, Lemme \ref{algebra-lemma-finite-flat-local}.
\item Si $f : A \to B$ est un morphisme d'anneaux quelconques,
$f$ est plat si et seulement si les homomorphismes induits
$A_{f^{-1}(\mathfrak q)} \to B_{\mathfrak q}$ sont plats pour tout
$\mathfrak q \in \Spec(B)$.
Voir Algèbre, Lemme \ref{algebra-lemma-flat-localization}.
\item Si $f : A \to B$ est un homomorphisme local d'anneaux locaux,
$f$ est plat si et seulement s'il est fidèlement plat.
Voir Algèbre, Lemme \ref{algebra-lemma-local-flat-ff}.
\item Un homomorphisme d'anneaux $A \to B$ est fidèlement plat si et seulement s'il est
plat et si l'application induite sur les spectres est surjective.
Voir Algèbre, Lemme \ref{algebra-lemma-ff-rings}.
\item Si $A$ est un anneau local noethérien, le complété
$A^\wedge$ est fidèlement plat sur $A$.
Voir Algèbre, Lemme \ref{algebra-lemma-completion-faithfully-flat}.
\item Soient $A$ un anneau local noethérien et $M$ un $A$-module.
Alors $M$ est plat sur $A$ si et seulement si $M \otimes_A A^\wedge$
est plat sur $A^\wedge$. (Combiner l'assertion précédente avec
Algèbre, Lemme \ref{algebra-lemma-flatness-descends}.)
\end{enumerate}
Avant de passer à la catégorie géométrique, nous présentons le théorème de
Grothendieck, qui fournit un procédé commode pour construire des modules
plats.

\begin{theorem}
\label{theorem-flatness-grothendieck}
Soient $A$, $B$ des anneaux locaux noethériens.
Soit $f : A \to B$ un homomorphisme local.
Si $M$ est un $B$-module fini qui est plat comme $A$-module,
et si $t \in \mathfrak m_B$ est un élément tel que la multiplication
par $t$ soit injective sur $M/\mathfrak m_AM$, alors $M/tM$ est également plat sur $A$.
\end{theorem}

\begin{proof}
Voir Algèbre, Lemme \ref{algebra-lemma-mod-injective}.
Voir aussi \cite[Section 20]{MatCA}.
\end{proof}

\begin{definition}
\label{definition-flat-schemes}
(Voir Morphismes, Définition \ref{morphisms-definition-flat}.)
Soit $f : X \to Y$ un morphisme de schémas.
Soit $\mathcal{F}$ un $\mathcal{O}_X$-module quasi-cohérent.
\begin{enumerate}
\item Soit $x \in X$. Nous disons que $\mathcal{F}$ est
{\it plat sur $Y$ en $x \in X$} si $\mathcal{F}_x$
est un $\mathcal{O}_{Y, f(x)}$-module plat.
On utilise l'homomorphisme $\mathcal{O}_{Y, f(x)} \to \mathcal{O}_{X, x}$ pour
considérer $\mathcal{F}_x$ comme un $\mathcal{O}_{Y, f(x)}$-module.
\item Soit $x \in X$. Nous disons que $f$ est {\it plat en $x \in X$}
si $\mathcal{O}_{Y, f(x)} \to \mathcal{O}_{X, x}$ est plat.
\item Nous disons que $f$ est {\it plat} s'il est plat en tout point de $X$.
\item Un morphisme $f : X \to Y$ plat et surjectif est parfois
dit {\it fidèlement plat}.
\end{enumerate}
\end{definition}

\noindent
Voici, une fois encore, une liste de résultats :
\begin{enumerate}
\item La propriété (pour un morphisme) d'être plat est, par définition,
locale pour la topologie de Zariski sur la source et la cible.
\item Les immersions ouvertes sont plates. (C'est clair, puisqu'elles induisent des isomorphismes
sur les anneaux locaux.)
\item Les morphismes plats sont stables par changement de base et par composition.
Morphismes, Lemmes \ref{morphisms-lemma-base-change-flat} et
\ref{morphisms-lemma-composition-flat}.
\item Si $f : X \to Y$ est plat, alors le foncteur image inverse
$\QCoh(\mathcal{O}_Y) \to \QCoh(\mathcal{O}_X)$ est exact.
Cela résulte immédiatement de l'examen des fibres.
\item Soit $f : X \to Y$ un morphisme de schémas, et supposons $Y$
quasi-compact et quasi-séparé. Dans ce cas,
si le foncteur $f^*$ est exact, alors $f$ est plat.
(Démonstration omise. Indication : utiliser
Propriétés, Lemme \ref{properties-lemma-extend-trivial} pour voir que
$Y$ possède « assez » de faisceaux d'idéaux et utiliser la caractérisation de la
platitude dans Algèbre, Lemme \ref{algebra-lemma-flat}.)
\end{enumerate}



\section{Propriétés topologiques des morphismes plats}
\label{section-topological-flat}

\noindent
Nous « rappelons » ci-dessous quelques propriétés d'ouverture dont jouissent les morphismes plats.

\begin{theorem}
\label{theorem-flat-open}
Soit $Y$ un schéma localement noethérien.
Soit $f : X \to Y$ un morphisme localement de type fini.
Soit $\mathcal{F}$ un $\mathcal{O}_X$-module cohérent.
L'ensemble des points de $X$ où $\mathcal{F}$ est plat sur $Y$ est ouvert.
En particulier, l'ensemble des points où $f$ est plat est ouvert dans $X$.
\end{theorem}

\begin{proof}
Voir Compléments sur les morphismes, Théorème \ref{more-morphisms-theorem-openness-flatness}.
\end{proof}

\begin{theorem}
\label{theorem-flat-map-open}
Soit $Y$ un schéma localement noethérien.
Soit $f : X \to Y$ un morphisme plat et localement de type fini.
Alors $f$ est (universellement) ouvert.
\end{theorem}

\begin{proof}
Voir Morphismes, Lemme \ref{morphisms-lemma-fppf-open}.
\end{proof}

\begin{theorem}
\label{theorem-flat-is-quotient}
Un morphisme fidèlement plat quasi-compact est une application quotient pour
la topologie de Zariski.
\end{theorem}

\begin{proof}
Voir Morphismes, Lemme \ref{morphisms-lemma-fpqc-quotient-topology}.
\end{proof}

\noindent
Une raison importante d'étudier les morphismes plats est qu'ils fournissent le cadre
adéquat pour saisir la notion de famille de schémas paramétrée par les
points d'un autre schéma. Naïvement, on pourrait penser que tout morphisme $f : X \to S$
doit être considéré comme une famille paramétrée par les points de $S$. Cependant,
sans hypothèse de platitude sur $f$, des phénomènes véritablement bizarres peuvent se produire dans
cette prétendue famille. Par exemple, rien ne garantit que la dimension relative
(la dimension des fibres) soit constante dans une famille. D'autres invariants numériques,
comme le polynôme de Hilbert, peuvent eux aussi varier d'une fibre à
l'autre. La platitude empêche de tels phénomènes et confère donc
une certaine « continuité » aux fibres.


\section{Morphismes \'etales}
\label{section-etale-morphisms}

\noindent
Dans cette section, nous allons définir les morphismes \'etales et démontrer plusieurs
propriétés importantes à leur sujet. La plus importante est sans aucun doute la
caractérisation fonctorielle présentée dans le Théorème \ref{theorem-formally-etale}.
Ensuite, nous examinerons aussi quelques propriétés des anneaux qui sont
insensibles à une extension \'etale (propriétés qui valent pour un anneau
si et seulement si elles valent pour toutes ses extensions \'etales), afin de motiver le
principe fondamental de la cohomologie \'etale --- les morphismes \'etales sont l'analogue algébrique des
isomorphismes locaux.

\medskip\noindent
Comme le titre le suggère, nous allons définir la classe des morphismes \'etales --- la
classe des morphismes (dont nous considérerons les familles surjectives comme des recouvrements)
dans la catégorie des schémas au-dessus d'un schéma de base $S$, afin de définir le
site \'etale $S_\etale$. Intuitivement, un morphisme \'etale est censé
exprimer l'idée d'un espace de revêtement et devrait donc être proche d'un
isomorphisme local. Si nous travaillons avec des variétés sur des corps algébriquement clos,
cette dernière assertion peut servir de définition, à condition de remplacer
« isomorphisme local » par « isomorphisme local formel » (isomorphisme après
complétion). On peut alors donner une définition sur un corps de base quelconque en demandant
que le changement de base à la clôture algébrique soit \'etale (au
sens précité). Mais, plutôt que de procéder par de telles constructions peu
esthétiques, nous adopterons une approche algébrique plus nette, quoique légèrement plus
abstraite.

\medskip\noindent
Nous définissons d'abord les « homomorphismes \'etales d'anneaux locaux » pour les anneaux
locaux noethériens. Nous ne pouvons pas employer simplement le terme « \'etale », puisqu'il
existe déjà une notion d'homomorphisme d'anneaux \'etale
(Algèbre, Section \ref{algebra-section-etale})
et qu'elle est différente.

\begin{definition}
\label{definition-etale-ring}
Soient $A$, $B$ des anneaux locaux noethériens.
Un homomorphisme local $f : A \to B$ est appelé
{\it homomorphisme \'etale d'anneaux locaux}
s'il est plat et non ramifié en tant qu'homomorphisme d'anneaux locaux
(voir Définition \ref{definition-unramified-rings}).
\end{definition}

\noindent
Il s'agit de la version locale de la définition d'un homomorphisme d'anneaux \'etale dans
Algèbre, Section \ref{algebra-section-etale}.
La définition précise
donnée dans cette section est celle d'un homomorphisme d'anneaux lisse de dimension
relative $0$. Il est démontré (dans
Algèbre, Lemme \ref{algebra-lemma-etale-standard-smooth})
qu'une $R$-algèbre \'etale $S$ possède toujours une présentation
$$
S = R[x_1, \ldots, x_n]/(f_1, \ldots, f_n)
$$
telle que
$$
g =
\det
\left(
\begin{matrix}
\partial f_1/\partial x_1 &
\partial f_2/\partial x_1 &
\ldots &
\partial f_n/\partial x_1 \\
\partial f_1/\partial x_2 &
\partial f_2/\partial x_2 &
\ldots &
\partial f_n/\partial x_2 \\
\ldots & \ldots & \ldots & \ldots \\
\partial f_1/\partial x_n &
\partial f_2/\partial x_n &
\ldots &
\partial f_n/\partial x_n
\end{matrix}
\right)
$$
ait pour image un élément inversible de $S$. Les deux lemmes suivants relient les deux
notions.

\begin{lemma}
\label{lemma-characterize-etale-Noetherian}
Soit $A \to B$ un homomorphisme de type fini, où $A$ est un anneau noethérien.
Soit $\mathfrak q$ un idéal premier de $B$ au-dessus de $\mathfrak p \subset A$.
Alors $A \to B$ est \'etale en $\mathfrak q$ si et seulement si
$A_{\mathfrak p} \to B_{\mathfrak q}$ est un homomorphisme \'etale
d'anneaux locaux.
\end{lemma}

\begin{proof}
Voir Algèbre, Lemmes \ref{algebra-lemma-etale} (platitude des homomorphismes \'etales),
\ref{algebra-lemma-etale-at-prime} (les homomorphismes \'etales sont non ramifiés)
et \ref{algebra-lemma-characterize-etale} (les homomorphismes plats et non ramifiés
sont \'etales).
\end{proof}

\begin{lemma}
\label{lemma-characterize-etale-completions}
Soient $A$, $B$ des anneaux locaux noethériens.
Soit $A \to B$ un homomorphisme local tel que $B$ soit essentiellement de
type fini sur $A$.
Les assertions suivantes sont équivalentes :
\begin{enumerate}
\item $A \to B$ est un homomorphisme \'etale d'anneaux locaux
\item $A^\wedge \to B^\wedge$ est un homomorphisme \'etale d'anneaux locaux, et
\item $A^\wedge \to B^\wedge$ est \'etale.
\end{enumerate}
De plus, dans ce cas $B^\wedge \cong (A^\wedge)^{\oplus n}$ comme
$A^\wedge$-modules pour un certain $n \geq 1$.
\end{lemma}

\begin{proof}
Pour établir les équivalences de (1), (2) et (3), puisque nous disposons des résultats
correspondants pour les homomorphismes d'anneaux non ramifiés
(Lemme \ref{lemma-characterize-unramified-completions}),
il suffit de démontrer que
$A \to B$ est plat si et seulement si $A^\wedge \to B^\wedge$ est plat.
Cela ressort de nos listes de propriétés des homomorphismes plats, puisque
les homomorphismes d'anneaux $A \to A^\wedge$ et $B \to B^\wedge$ sont fidèlement plats.
Pour la dernière assertion, Lemme \ref{lemma-unramified-completions}
montre que $B^\wedge$ est un $A^\wedge$-module fini et plat.
Il est donc fini et libre d'après notre liste
des propriétés des modules plats de la Section \ref{section-flat-morphisms}.
\end{proof}

\noindent
L'entier $n$ qui intervient dans le lemme ci-dessus
n'est autre que le degré
$[\kappa(\mathfrak m_B) : \kappa(\mathfrak m_A)]$ de l'extension des corps
résiduels. En particulier, si $\kappa(\mathfrak m_A)$
est séparablement clos, nous voyons que $A^\wedge \to B^\wedge$
est un isomorphisme, ce qui confirme nos assertions précédentes.

\begin{definition}
\label{definition-etale-schemes-1}
(Voir Morphismes, Définition \ref{morphisms-definition-etale}.)
Soit $Y$ un schéma localement noethérien.
Soit $f : X \to Y$ un morphisme de schémas localement de type fini.
\begin{enumerate}
\item Soit $x \in X$. Nous disons que $f$ est {\it \'etale en $x \in X$} si
$\mathcal{O}_{Y, f(x)} \to \mathcal{O}_{X, x}$ est un
homomorphisme \'etale d'anneaux locaux.
\item Le morphisme est dit {\it \'etale} s'il est \'etale en chacun de ses
points.
\end{enumerate}
\end{definition}

\noindent
Démontrons que cette définition coïncide avec celle du
chapitre sur les morphismes de schémas. Cela garantit en particulier que
l'ensemble des points où un morphisme est \'etale est ouvert.

\begin{lemma}
\label{lemma-etale-definition}
Soit $Y$ un schéma localement noethérien.
Soit $f : X \to Y$ localement de type fini.
Soit $x \in X$. Le morphisme $f$ est \'etale en $x$ au
sens de la Définition \ref{definition-etale-schemes-1}
si et seulement s'il est \'etale en $x$ au
sens de Morphismes, Définition \ref{morphisms-definition-etale}.
\end{lemma}

\begin{proof}
Cela résulte du Lemme \ref{lemma-characterize-etale-Noetherian}
et des définitions.
\end{proof}

\noindent
Voici quelques résultats sur les morphismes \'etales.
Les formulations données dans cette liste ne s'appliquent qu'aux
morphismes localement de type fini entre schémas localement noethériens.
Dans chaque cas, nous donnons une référence au résultat général tel qu'il a été
démontré plus tôt dans le projet, mais on peut, dans certains cas,
démontrer plus facilement le résultat dans le cas noethérien.
Voici la liste :
\begin{enumerate}
\item Un morphisme \'etale est non ramifié. (Cela ressort de nos définitions.)
\item La propriété d'être \'etale est locale sur la source et la cible pour la
topologie de Zariski.
\item Les morphismes \'etales sont stables par changement de base et par composition.
Voir Morphismes, Lemmes \ref{morphisms-lemma-base-change-etale}
et \ref{morphisms-lemma-composition-etale}.
\item Les morphismes \'etales de schémas sont localement quasi-finis,
et les morphismes \'etales quasi-compacts sont quasi-finis. (Cela est
vrai parce que c'est le cas pour les morphismes non ramifiés, comme vu précédemment.)
\item Les morphismes \'etales sont de dimension relative $0$. Voir
Morphismes, Définition \ref{morphisms-definition-relative-dimension-d}
et
Morphismes, Lemme \ref{morphisms-lemma-locally-quasi-finite-rel-dimension-0}.
\item Un morphisme est \'etale si et seulement s'il est plat et si
toutes ses fibres sont \'etales. Voir
Morphismes, Lemme \ref{morphisms-lemma-etale-flat-etale-fibres}.
\item Les morphismes \'etales sont ouverts. C'est vrai parce qu'un morphisme \'etale
est plat, d'après le Théorème \ref{theorem-flat-map-open}.
\item Soient $X$ et $Y$ \'etales sur un schéma de base $S$.
Tout $S$-morphisme de $X$ vers $Y$ est \'etale.
Voir Morphismes, Lemme \ref{morphisms-lemma-etale-permanence}.
\end{enumerate}






\section{Le théorème de structure}
\label{section-structure-etale-map}

\noindent
Nous présentons un théorème qui décrit la structure locale des morphismes \'etales
et non ramifiés. Outre son importance propre évidente,
ce théorème nous permet aussi de passer à une autre
définition des morphismes \'etales qui rend mieux compte de l'intuition géométrique
que celle que nous avons employée jusqu'ici.

\medskip\noindent
Pour l'énoncer, nous avons besoin de la notion d'{\it homomorphisme d'anneaux \'etale standard}, voir
Algèbre, Définition \ref{algebra-definition-standard-etale}.
Plus précisément, supposons que $R$ soit un anneau et que $f, g \in R[t]$ soient des polynômes
tels que
\begin{enumerate}
\item[(a)] $f$ soit un polynôme unitaire, et
\item[(b)] $f' = \text{d}f/\text{d}t$ soit inversible dans le localisé
$R[t]_g/(f)$.
\end{enumerate}
Alors l'application
$$
R \longrightarrow R[t]_g/(f) = R[t, 1/g]/(f)
$$
est une algèbre \'etale standard, et toute algèbre \'etale standard est isomorphe
à l'une d'elles. C'est un exercice agréable de démontrer qu'un tel homomorphisme d'anneaux
est plat et non ramifié, donc \'etale (comme prévu, bien entendu).
Un cas particulier d'homomorphisme d'anneaux \'etale standard est tout homomorphisme d'anneaux
$$
R \longrightarrow R[t]_{f'}/(f) = R[t, 1/f']/(f)
$$
où $f$ est un polynôme unitaire, et toute algèbre \'etale standard est (isomorphe à)
un localisé principal de l'une d'elles.

\begin{theorem}
\label{theorem-structure-etale}
Soit $f : A \to B$ un homomorphisme \'etale d'anneaux locaux.
Alors il existe $f, g \in A[t]$ tels que
\begin{enumerate}
\item $B' = A[t]_g/(f)$ soit une algèbre \'etale standard --- voir (a) et (b) ci-dessus, et
\item $B$ soit isomorphe au localisé de $B'$ en un idéal premier.
\end{enumerate}
\end{theorem}

\begin{proof}
Écrivons $B = B'_{\mathfrak q}$ pour une certaine $A$-algèbre de type fini $B'$
(ce qui est possible puisque $B$ est essentiellement de type fini sur $A$).
D'après le Lemme \ref{lemma-characterize-etale-Noetherian},
nous voyons que $A \to B'$ est \'etale en $\mathfrak q$.
Nous pouvons donc appliquer
Algèbre, Proposition \ref{algebra-proposition-etale-locally-standard}
pour voir qu'un localisé principal de $B'$ est \'etale standard.
\end{proof}

\noindent
Voici la version pour les homomorphismes non ramifiés d'anneaux locaux.

\begin{theorem}
\label{theorem-structure-unramified}
Soit $f : A \to B$ un homomorphisme non ramifié d'anneaux locaux.
Alors il existe $f, g \in A[t]$ tels que
\begin{enumerate}
\item $B' = A[t]_g/(f)$ soit une algèbre \'etale standard --- voir (a) et (b) ci-dessus, et
\item $B$ soit isomorphe à un quotient d'un localisé de $B'$ en un idéal premier.
\end{enumerate}
\end{theorem}

\begin{proof}
Écrivons $B = B'_{\mathfrak q}$ pour une certaine $A$-algèbre de type fini $B'$
(ce qui est possible puisque $B$ est essentiellement de type fini sur $A$).
D'après le Lemme \ref{lemma-characterize-unramified-Noetherian},
nous voyons que $A \to B'$ est non ramifié en $\mathfrak q$.
Nous pouvons donc appliquer
Algèbre, Proposition \ref{algebra-proposition-unramified-locally-standard}
pour voir qu'un localisé principal de $B'$ est un quotient d'une
$A$-algèbre \'etale standard.
\end{proof}

\noindent
Des arguments usuels de relèvement donnent alors l'énoncé géométrique
suivant, qui nous sera d'une utilité essentielle.

\begin{theorem}
\label{theorem-geometric-structure}
Soit $\varphi : X \to Y$ un morphisme de schémas. Soit $x \in X$.
Soit $V \subset Y$ un voisinage ouvert affine de $\varphi(x)$.
Si $\varphi$ est \'etale en $x$, alors il existe un ouvert affine
$U \subset X$ avec $x \in U$ et $\varphi(U) \subset V$
tel que l'on ait le diagramme suivant :
$$
\xymatrix{
X \ar[d] & U \ar[l] \ar[d] \ar[r]_-j & \Spec(R[t]_{f'}/(f)) \ar[d] \\
Y & V \ar[l] \ar@{=}[r] & \Spec(R)
}
$$
où $j$ est une immersion ouverte et $f \in R[t]$ est un polynôme unitaire.
\end{theorem}

\begin{proof}
Cet énoncé équivaut à celui de
Morphismes, Lemme \ref{morphisms-lemma-etale-locally-standard-etale},
bien que les énoncés diffèrent légèrement.
Voir aussi Variétés, Lemme \ref{varieties-lemma-geometric-structure-unramified}
pour une variante relative aux morphismes non ramifiés.
\end{proof}


\section{Morphismes \'etales et lisses}
\label{section-etale-smooth}

\noindent
Un morphisme \'etale est lisse de dimension relative zéro.
La projection $\mathbf{A}^n_S \to S$ est un exemple standard
de morphisme lisse de dimension relative $n$.
Il se trouve que, localement pour la topologie \'etale, tout morphisme lisse
est de cette forme. Voici l'énoncé précis.

\begin{theorem}
\label{theorem-smooth-etale-over-n-space}
Soit $\varphi : X \to Y$ un morphisme de schémas.
Soit $x \in X$.
Si $\varphi$ est lisse en $x$, alors
il existe un entier $n \geq 0$ et des ouverts affines
$V \subset Y$ et $U \subset X$ avec $x \in U$ et $\varphi(U) \subset V$
de sorte qu'il existe un diagramme commutatif
$$
\xymatrix{
X \ar[d] & U \ar[l] \ar[d] \ar[r]_-\pi &
\mathbf{A}^n_R \ar[d] \ar@{=}[r] &  \Spec(R[x_1, \ldots, x_n]) \ar[dl] \\
Y & V \ar[l] \ar@{=}[r] & \Spec(R)
}
$$
où $\pi$ est \'etale.
\end{theorem}

\begin{proof}
Voir
Morphismes, Lemme \ref{morphisms-lemma-smooth-etale-over-affine-space}.
\end{proof}




\section{Propriétés topologiques des morphismes \'etales}
\label{section-topological-etale}

\noindent
Nous présentons quelques propriétés topologiques des morphismes \'etales et
non ramifiés. Tout d'abord, nous donnons ce que Grothendieck
appelle la {\it propriété fondamentale des morphismes \'etales}, voir
\cite[Expos\'e I.5]{SGA1}.

\begin{theorem}
\label{theorem-etale-radicial-open}
Soit $f : X \to Y$ un morphisme de schémas.
Les assertions suivantes sont équivalentes :
\begin{enumerate}
\item $f$ est une immersion ouverte,
\item $f$ est universellement injectif et \'etale, et
\item $f$ est un monomorphisme plat, localement de présentation finie.
\end{enumerate}
\end{theorem}

\begin{proof}
Une immersion ouverte est universellement injective,
car tout changement de base d'une immersion ouverte
est une immersion ouverte. De plus, elle est \'etale d'après
Morphismes, Lemme \ref{morphisms-lemma-open-immersion-etale}.
Ainsi (1) implique (2).

\medskip\noindent
Supposons que $f$ soit universellement injectif et \'etale.
Comme $f$ est \'etale, il est plat et localement de présentation finie, voir
Morphismes, Lemmes \ref{morphisms-lemma-etale-flat} et
\ref{morphisms-lemma-etale-locally-finite-presentation}.
D'après
Lemme \ref{lemma-universally-injective-unramified},
nous voyons que $f$ est un monomorphisme. Ainsi (2) implique (3).

\medskip\noindent
Supposons que $f$ soit plat, localement de présentation finie et un monomorphisme.
Alors $f$ est ouvert, voir
Morphismes, Lemme \ref{morphisms-lemma-fppf-open}.
Nous pouvons donc remplacer $Y$ par $f(X)$ et supposer que $f$ est
surjectif. Alors $f$ est ouvert et bijectif, donc c'est un homéomorphisme.
Par conséquent, $f$ est quasi-compact. Dès lors,
Descente, lemme
\ref{descent-lemma-flat-surjective-quasi-compact-monomorphism-isomorphism}
montre que $f$ est un isomorphisme, ce qui conclut.
\end{proof}

\noindent
Voici un autre résultat de même nature.

\begin{lemma}
\label{lemma-finite-etale-one-point}
Soit $\pi : X \to S$ un morphisme de schémas. Soit $s \in S$.
Supposons que
\begin{enumerate}
\item $\pi$ soit fini ;
\item $\pi$ soit \'etale ;
\item $\pi^{-1}(\{s\}) = \{x\}$ ; et
\item $\kappa(s) \subset \kappa(x)$ soit purement
inséparable\footnote{Compte tenu de la condition (2)
cela équivaut à $\kappa(s) = \kappa(x)$.}.
\end{enumerate}
Alors il existe un voisinage ouvert $U$ de $s$ tel que
$\pi|_{\pi^{-1}(U)} : \pi^{-1}(U) \to U$ soit un isomorphisme.
\end{lemma}

\begin{proof}
D'après le
lemme \ref{lemma-finite-unramified-one-point},
il existe un voisinage ouvert $U$ de $s$ tel que
$\pi|_{\pi^{-1}(U)} : \pi^{-1}(U) \to U$ soit une immersion fermée.
Mais un morphisme qui est à la fois \'etale et une immersion fermée est une
immersion ouverte (par exemple d'après le
théorème \ref{theorem-etale-radicial-open}).
En rétrécissant $U$, on obtient donc un isomorphisme.
\end{proof}

\begin{lemma}
\label{lemma-relative-frobenius-etale}
Soit $U \to X$ un morphisme \'etale de schémas
où $X$ est un schéma de caractéristique $p$.
Alors le Frobenius relatif $F_{U/X} : U \to U \times_{X, F_X} X$
est un isomorphisme.
\end{lemma}

\begin{proof}
Le morphisme $F_{U/X}$ est un homéomorphisme universel d'après
Variétés, lemme \ref{varieties-lemma-relative-frobenius}.
Le morphisme $F_{U/X}$ est \'etale en tant que
morphisme entre des schémas \'etales sur $X$
(Morphismes, lemme \ref{morphisms-lemma-etale-permanence}).
Par conséquent, $F_{U/X}$ est un isomorphisme d'après le
théorème \ref{theorem-etale-radicial-open}.
\end{proof}






\section{Invariance topologique de la topologie \'etale}
\label{section-topological-invariance}

\noindent
Nous présentons maintenant un théorème absolument crucial qui affirme, en gros,
que le caractère \'etale est une propriété topologique.

\begin{theorem}
\label{theorem-etale-topological}
Soient $X$ et $Y$ deux schémas sur un schéma de base $S$. Soit $S_0$ un sous-schéma fermé
de $S$ ayant le même espace topologique sous-jacent
(par exemple si le faisceau d'idéaux de $S_0$ dans $S$ est de carré nul).
Notons $X_0$ (resp.\ $Y_0$) le changement de base $S_0 \times_S X$
(resp.\ $S_0 \times_S Y$).
Si $X$ est \'etale sur $S$, alors l'application
$$
\Mor_S(Y, X) \longrightarrow \Mor_{S_0}(Y_0, X_0)
$$
est bijective.
\end{theorem}

\begin{proof}
Après le changement de base par $Y \to S$, nous pouvons supposer que $Y = S$.
Dans ce cas, le théorème affirme que tout $S$-morphisme $\sigma_0 : S_0 \to X$
se factorise en fait de manière unique par une section $S \to X$ du
morphisme structural \'etale $f : X \to S$.

\medskip\noindent
Unicité. Supposons que nous ayons deux sections $\sigma, \sigma'$
par lesquelles $\sigma_0$ se factorise. Comme $X \to S$ est \'etale,
nous voyons que $\Delta : X \to X \times_S X$ est une immersion ouverte
(Morphismes, lemme \ref{morphisms-lemma-diagonal-unramified-morphism}).
Le morphisme $(\sigma, \sigma') : S \to X \times_S X$ se factorise
par cet ouvert, car, pour tout $s \in S$, nous avons
$(\sigma, \sigma')(s) = (\sigma_0(s), \sigma_0(s))$. Ainsi,
$\sigma = \sigma'$.

\medskip\noindent
Pour démontrer l'existence, nous nous ramenons d'abord au cas affine
(nous suggérons au lecteur de sauter cette étape).
Soit $X = \bigcup X_i$ un recouvrement par des ouverts affines tel
que chaque $X_i$ s'envoie dans un ouvert affine $S_i$ de $S$.
Pour tout $s \in S$, nous pouvons choisir $i$ tel que
$\sigma_0(s) \in X_i$.
Choisissons un voisinage ouvert affine $U \subset S_i$ de $s$
tel que $\sigma_0(U_0) \subset X_{i, 0}$. Remarquons que
$X' = X_i \times_S U = X_i \times_{S_i} U$ est affine.
Si l'on peut relever $\sigma_0|_{U_0} : U_0 \to X'_0$ en
un morphisme $U \to X'$, alors, par unicité, ces relèvements locaux se recollent
en un morphisme global $S \to X$. Nous pouvons donc supposer $S$ et
$X$ affines.

\medskip\noindent
Existence lorsque $S$ et $X$ sont affines. Écrivons $S = \Spec(A)$
et $X = \Spec(B)$. Alors $A \to B$ est \'etale et, en particulier,
lisse (de dimension relative $0$). Puisque $|S_0| = |S|$, nous voyons
que $S_0 = \Spec(A/I)$, où $I \subset A$ est localement nilpotent.
L'existence résulte donc de
Algèbre, lemme \ref{algebra-lemma-smooth-strong-lift}.
\end{proof}

\noindent
La démonstration du théorème précédent nous donne aussi une implication de la
caractérisation fonctorielle promise des morphismes \'etales. Le
théorème suivant sera renforcé dans
Cohomologie \'etale,
théorème \ref{etale-cohomology-theorem-topological-invariance}.

\begin{theorem}[Une équivalence remarquable de cat\'egories]
\label{theorem-remarkable-equivalence}
\begin{reference}
\cite[IV, Theorem 18.1.2]{EGA}
\end{reference}
Soit $S$ un schéma.
Soit $S_0 \subset S$ un sous-schéma fermé ayant le même espace
topologique sous-jacent (par exemple si le faisceau d'idéaux de $S_0$ dans $S$
est de carré nul). Le foncteur
$$
X \longmapsto X_0 = S_0 \times_S X
$$
définit une équivalence de catégories
$$
\{
\text{schémas }X\text{ \'etales sur }S
\}
\leftrightarrow
\{
\text{schémas }X_0\text{ \'etales sur }S_0
\}
$$
\end{theorem}

\begin{proof}
Le théorème \ref{theorem-etale-topological}
montre que ce foncteur est pleinement fidèle.
Il reste à démontrer qu'il est essentiellement surjectif.
Soit $Y \to S_0$ un morphisme \'etale de schémas.

\medskip\noindent
Supposons le résultat vrai lorsque $S$ et $Y$ sont affines.
Dans ce cas, choisissons un recouvrement par des ouverts affines
$Y = \bigcup V_j$ tel que chaque $V_j$ s'envoie
dans un ouvert affine de $S$. Par hypothèse (cas affine), nous pouvons
trouver des morphismes \'etales $W_j \to S$ tels que $W_{j, 0} \cong V_j$
(comme schémas sur $S_0$). Soit $W_{j, j'} \subset W_j$
le sous-schéma ouvert dont l'espace topologique sous-jacent
correspond à $V_j \cap V_{j'}$. Puisque nous avons des isomorphismes
$$
W_{j, j', 0} \cong V_j \cap V_{j'} \cong W_{j', j, 0}
$$
de schémas sur $S_0$, la pleine fidélité montre que nous
obtenons des isomorphismes
$\theta_{j, j'} : W_{j, j'} \to W_{j', j}$ de schémas sur $S$.
Nous omettons de vérifier que ces isomorphismes satisfont à la
condition de cocycle de Schémas, section \ref{schemes-section-glueing-schemes}.
En appliquant Schémas, lemme \ref{schemes-lemma-glue-schemes}
nous obtenons un schéma $X \to S$ en
recollant les schémas $W_j$ au moyen des identifications $\theta_{j, j'}$.
Il est clair, par construction, que $X \to S$ est \'etale et que $X_0 \cong Y$.

\medskip\noindent
Il suffit donc de démontrer le lemme lorsque $S$ et $Y$ sont affines.
Écrivons $S = \Spec(R)$ et $S_0 = \Spec(R/I)$, où $I$ est localement nilpotent.
D'après Algèbre, lemme \ref{algebra-lemma-etale-standard-smooth}, nous savons que
$Y$ est le spectre d'un anneau $\overline{A}$ avec
$$
\overline{A} = (R/I)[x_1, \ldots, x_n]/(\overline{f}_1, \ldots, \overline{f}_n)
$$
tel que
$$
\overline{g} =
\det
\left(
\begin{matrix}
\partial \overline{f}_1/\partial x_1 &
\partial \overline{f}_2/\partial x_1 &
\ldots &
\partial \overline{f}_n/\partial x_1 \\
\partial \overline{f}_1/\partial x_2 &
\partial \overline{f}_2/\partial x_2 &
\ldots &
\partial \overline{f}_n/\partial x_2 \\
\ldots & \ldots & \ldots & \ldots \\
\partial \overline{f}_1/\partial x_n &
\partial \overline{f}_2/\partial x_n &
\ldots &
\partial \overline{f}_n/\partial x_n
\end{matrix}
\right)
$$
a pour image un élément inversible dans $\overline{A}$. Choisissons des relèvements quelconques
$f_i \in R[x_1, \ldots, x_n]$. Posons
$$
A = R[x_1, \ldots, x_n]/(f_1, \ldots, f_n)
$$
Puisque $I$ est localement nilpotent, l'idéal $IA$ est localement nilpotent
(Algèbre, lemme \ref{algebra-lemma-locally-nilpotent}).
Remarquons que $\overline{A} = A/IA$.
Il s'ensuit que le déterminant de la matrice des dérivées partielles des
$f_i$ est inversible dans l'algèbre $A$, d'après
Algèbre, lemme \ref{algebra-lemma-locally-nilpotent-unit}.
Par conséquent, $R \to A$ est \'etale, ce qui achève la démonstration.
\end{proof}



\section{La caractérisation fonctorielle}
\label{section-functorial-etale}

\noindent
Nous présentons enfin la caractérisation fonctorielle annoncée.
Il y a donc quatre façons d'envisager les morphismes \'etales de schémas :
\begin{enumerate}
\item comme des morphismes lisses de dimension relative $0$,
\item comme des morphismes localement de présentation finie, plats et non ramifiés,
\item au moyen du théorème de structure, et
\item au moyen de la caractérisation fonctorielle.
\end{enumerate}

\begin{theorem}
\label{theorem-formally-etale}
Soit $f : X \to S$ un morphisme localement de présentation finie.
Les assertions suivantes sont équivalentes
\begin{enumerate}
\item $f$ est \'etale,
\item pour tout $S$-schéma affine $Y$ et tout sous-schéma fermé $Y_0 \subset Y$
défini par un idéal de carré nul, l'application naturelle
$$
\Mor_S(Y, X) \longrightarrow \Mor_S(Y_0, X)
$$
est bijective.
\end{enumerate}
\end{theorem}

\begin{proof}
C'est
Compléments sur les morphismes, lemme \ref{more-morphisms-lemma-etale-formally-etale}.
\end{proof}

\noindent
Cette caractérisation affirme que les solutions des équations définissant $X$ se
relèvent de façon unique à travers les épaississements nilpotents.



\section{Structure locale \'etale des morphismes non ramifiés}
\label{section-unramified-etale-local}

\noindent
Au chapitre
Compléments sur les morphismes, section \ref{more-morphisms-section-etale-localization},
le lecteur trouvera quelques résultats sur la structure locale \'etale des
morphismes quasi-finis. Dans cette section, nous voulons les combiner
avec les propriétés topologiques des morphismes non ramifiés rencontrées
dans ce chapitre. La configuration générale à garder à l'esprit est
$$
\xymatrix{
V \ar[r] \ar[dr] & X_U \ar[d] \ar[r] & X \ar[d]^f \\
& U \ar[r] & S
}
$$
voir
Compléments sur les morphismes, équation (\ref{more-morphisms-equation-basic-diagram}).
Commençons par un cas très général.

\begin{lemma}
\label{lemma-unramified-etale-local}
Soit $f : X \to S$ un morphisme de schémas.
Soient $x_1, \ldots, x_n \in X$ des points ayant même image $s$ dans $S$.
Supposons que $f$ soit non ramifié en chacun des $x_i$.
Il existe alors un voisinage \'etale $(U, u) \to (S, s)$
et des ouverts $V_{i, j} \subset X_U$, $i = 1, \ldots, n$, $j = 1, \ldots, m_i$,
tels que
\begin{enumerate}
\item $V_{i, j} \to U$ est une immersion fermée dont l'image contient $u$,
\item $u$ n'appartient pas à l'image de $V_{i, j} \cap V_{i', j'}$, sauf si
$i = i'$ et $j = j'$, et
\item tout point de $(X_U)_u$ s'envoyant sur $x_i$ appartient à l'un des $V_{i, j}$.
\end{enumerate}
\end{lemma}

\begin{proof}
D'après
Morphismes, définition \ref{morphisms-definition-unramified},
il existe, pour chaque $x_i$, un voisinage ouvert qui est localement de type
fini sur $S$. En remplaçant $X$ par un voisinage ouvert de $\{x_1, \ldots, x_n\}$,
nous pouvons supposer que $f$ est localement de type fini. Appliquons
Compléments sur les morphismes, lemme
\ref{more-morphisms-lemma-etale-makes-quasi-finite-finite-multiple-points-var}
pour obtenir le voisinage \'etale $(U, u)$ et les ouverts $V_{i, j}$ finis sur
$U$. D'après
le lemme \ref{lemma-finite-unramified-one-point},
quitte à rétrécir $U$, le morphisme $V_{i, j} \to U$ est une immersion
fermée.
\end{proof}

\begin{lemma}
\label{lemma-unramified-etale-local-technical}
Soit $f : X \to S$ un morphisme de schémas.
Soient $x_1, \ldots, x_n \in X$ des points ayant même image $s$ dans $S$.
Supposons que $f$ soit séparé et que $f$ soit non ramifié en chacun des $x_i$.
Il existe alors un voisinage \'etale $(U, u) \to (S, s)$
et une décomposition en somme disjointe
$$
X_U =
W \amalg \coprod\nolimits_{i, j} V_{i, j}
$$
telle que
\begin{enumerate}
\item $V_{i, j} \to U$ est une immersion fermée dont l'image contient $u$,
\item la fibre $W_u$ ne contient aucun point s'envoyant sur l'un des $x_i$.
\end{enumerate}
En particulier, si $f^{-1}(\{s\}) = \{x_1, \ldots, x_n\}$, alors
la fibre $W_u$ est vide.
\end{lemma}

\begin{proof}
Appliquons
le lemme \ref{lemma-unramified-etale-local}.
Nous pouvons supposer que $U$ est affine, de sorte que $X_U$ est séparé.
Alors $V_{i, j} \to X_U$ est une application fermée ; voir
Morphismes, lemme \ref{morphisms-lemma-image-proper-scheme-closed}.
Supposons que $(i, j) \not = (i', j')$.
Alors $V_{i, j} \cap V_{i', j'}$ est fermé dans $V_{i, j}$ et
son image dans $U$ ne contient pas $u$.
Par conséquent, quitte à rétrécir $U$, nous pouvons supposer que
$V_{i, j} \cap V_{i', j'} = \emptyset$. De plus, $\bigcup V_{i, j}$ est
un sous-schéma ouvert et fermé de $X_U$ et admet donc un
complémentaire ouvert et fermé $W$. Cela achève la démonstration.
\end{proof}

\noindent
Le lemme suivant est en un certain sens beaucoup plus faible que le précédent,
mais il peut être utile de l'énoncer explicitement ici. Il affirme qu'un morphisme fini
non ramifié est, localement pour la topologie \'etale sur la base, une immersion fermée.

\begin{lemma}
\label{lemma-finite-unramified-etale-local}
Soit $f : X \to S$ un morphisme fini non ramifié de schémas.
Soit $s \in S$.
Il existe un voisinage \'etale $(U, u) \to (S, s)$
et une décomposition en somme disjointe finie
$$
X_U = \coprod\nolimits_j V_j
$$
telle que chaque $V_j \to U$ soit une immersion fermée.
\end{lemma}

\begin{proof}
Puisque $X \to S$ est fini, la fibre au-dessus de $s$ est un ensemble fini
$\{x_1, \ldots, x_n\}$ de points de $X$. Appliquons le
lemme \ref{lemma-unramified-etale-local-technical}
à cet ensemble (un morphisme fini est séparé ; voir
Morphismes, section \ref{morphisms-section-integral}).
L'image de $W$ dans $U$ est un sous-ensemble fermé
(car $X_U \to U$ est fini, donc propre) qui ne
contient pas $u$. Après l'avoir retiré de $U$, nous voyons que $W = \emptyset$
comme voulu.
\end{proof}




\section{Structure locale \'etale des morphismes \'etales}
\label{section-etale-local-etale}

\noindent
C'est un peu naïf, mais peut-être utile pour se forger une intuition des morphismes \'etales.
Nous recopions simplement les résultats de la
section \ref{section-unramified-etale-local}
et remplaçons ``immersion fermée'' par ``isomorphisme''.

\begin{lemma}
\label{lemma-etale-etale-local}
Soit $f : X \to S$ un morphisme de schémas.
Soient $x_1, \ldots, x_n \in X$ des points ayant la même image $s$ dans $S$.
Supposons que $f$ soit \'etale en chacun des $x_i$.
Alors il existe un voisinage \'etale $(U, u) \to (S, s)$
et des ouverts $V_{i, j} \subset X_U$, $i = 1, \ldots, n$, $j = 1, \ldots, m_i$
tels que
\begin{enumerate}
\item $V_{i, j} \to U$ soit un isomorphisme ;
\item $u$ n'appartienne pas à l'image de $V_{i, j} \cap V_{i', j'}$, sauf si
$i = i'$ et $j = j'$ ; et
\item tout point de $(X_U)_u$ qui s'envoie sur $x_i$ appartienne à un certain $V_{i, j}$.
\end{enumerate}
\end{lemma}

\begin{proof}
Un morphisme \'etale est non ramifié ; nous pouvons donc appliquer le
lemme \ref{lemma-unramified-etale-local}.
Or $V_{i, j} \to U$ est une immersion fermée et est \'etale.
C'est donc une immersion ouverte, par exemple d'après le
théorème \ref{theorem-etale-radicial-open}.
Remplaçons $U$ par l'intersection des images des morphismes $V_{i, j} \to U$
pour obtenir le lemme.
\end{proof}

\begin{lemma}
\label{lemma-etale-etale-local-technical}
Soit $f : X \to S$ un morphisme de schémas.
Soient $x_1, \ldots, x_n \in X$ des points ayant la même image $s$ dans $S$.
Supposons $f$ séparé et $f$ \'etale en chacun des $x_i$.
Alors il existe un voisinage \'etale $(U, u) \to (S, s)$
et une décomposition en somme disjointe finie
$$
X_U =
W \amalg \coprod\nolimits_{i, j} V_{i, j}
$$
de schémas tels que
\begin{enumerate}
\item $V_{i, j} \to U$ soit un isomorphisme ;
\item la fibre $W_u$ ne contienne aucun point s'envoyant sur l'un des $x_i$.
\end{enumerate}
En particulier, si $f^{-1}(\{s\}) = \{x_1, \ldots, x_n\}$, alors
la fibre $W_u$ est vide.
\end{lemma}

\begin{proof}
Un morphisme \'etale est non ramifié ; nous pouvons donc appliquer le
lemme \ref{lemma-unramified-etale-local-technical}.
Comme dans la démonstration du
lemme \ref{lemma-etale-etale-local},
les morphismes $V_{i, j} \to U$ sont des immersions ouvertes et
l'on conclut en remplaçant $U$ par l'intersection de leurs
images.
\end{proof}

\noindent
Le lemme suivant est en un certain sens beaucoup plus faible que le précédent,
mais il peut être utile de l'énoncer explicitement ici. Il affirme qu'un morphisme fini
\'etale est, localement pour la topologie \'etale sur la base, un
``espace de revêtement topologique'', c'est-à-dire un produit fini de copies de la base.

\begin{lemma}
\label{lemma-finite-etale-etale-local}
Soit $f : X \to S$ un morphisme fini \'etale de schémas.
Soit $s \in S$. Il existe un voisinage \'etale $(U, u) \to (S, s)$
et une décomposition en somme disjointe finie
$$
X_U = \coprod\nolimits_j V_j
$$
de schémas telle que chaque $V_j \to U$ soit un isomorphisme.
\end{lemma}

\begin{proof}
Un morphisme \'etale est non ramifié ; nous pouvons donc appliquer le
lemme \ref{lemma-finite-unramified-etale-local}.
Comme dans la démonstration du
lemme \ref{lemma-etale-etale-local},
nous voyons que $V_{i, j} \to U$ est une immersion ouverte et l'on conclut
en remplaçant $U$ par l'intersection de leurs images.
\end{proof}




\section{Propriétés de permanence}
\label{section-properties-permanence}

\noindent
Dans ce qui suit, nous présentons quelques propriétés de ``permanence''
des homomorphismes \'etales d'anneaux locaux noethériens
(au sens de la définition \ref{definition-etale-ring}). Voir
Compléments d'algèbre, sections \ref{more-algebra-section-permanence-completion} et
\ref{more-algebra-section-permanence-henselization}
pour les résultats analogues concernant la complétion et
l'hensélisation d'un anneau local noethérien.

\begin{lemma}
\label{lemma-etale-dimension}
Soient $A$ et $B$ des anneaux locaux noethériens.
Soit $A \to B$ un homomorphisme local \'etale.
Alors $\dim(A) = \dim(B)$.
\end{lemma}

\begin{proof}
Voir par exemple
Algèbre, lemme \ref{algebra-lemma-dimension-base-fibre-equals-total}.
\end{proof}

\begin{proposition}
\label{proposition-etale-depth}
Soient $A$ et $B$ des anneaux locaux noethériens.
Soit $f : A \to B$ un homomorphisme local \'etale.
Alors $\text{profondeur}(A) = \text{profondeur}(B)$
\end{proposition}

\begin{proof}
Voir Algèbre, lemme \ref{algebra-lemma-apply-grothendieck}.
\end{proof}

\begin{proposition}
\label{proposition-etale-CM}
\begin{slogan}
La propriété d'être de Cohen-Macaulay monte et descend le long des homomorphismes \'etales.
\end{slogan}
Soient $A$ et $B$ des anneaux locaux noethériens.
Soit $f : A \to B$ un homomorphisme local \'etale.
Alors $A$ est de Cohen-Macaulay si et seulement si $B$ l'est.
\end{proposition}

\begin{proof}
Un anneau local $A$ est de Cohen-Macaulay si et seulement si $\dim(A) = \text{profondeur}(A)$.
Puisque ces deux invariants sont préservés par une extension \'etale,
l'assertion en résulte.
\end{proof}

\begin{proposition}
\label{proposition-etale-regular}
Soient $A$, $B$ des anneaux locaux noethériens.
Soit $f : A \to B$ un homomorphisme local \'etale.
Alors $A$ est régulier si et seulement si $B$ l'est.
\end{proposition}

\begin{proof}
Si $B$ est régulier, alors $A$ est régulier d'après
Algèbre, lemme \ref{algebra-lemma-flat-under-regular}.
Supposons $A$ régulier. Soit $\mathfrak m$ l'idéal maximal
de $A$. Alors $\dim_{\kappa(\mathfrak m)} \mathfrak m/\mathfrak m^2 =
\dim(A) = \dim(B)$ (voir le lemme \ref{lemma-etale-dimension}).
D'autre part, $\mathfrak mB$ est l'idéal maximal de
$B$ et par conséquent $\mathfrak m_B/\mathfrak m_B = \mathfrak mB/\mathfrak m^2B$
est engendré par au plus $\dim(B)$ éléments. Ainsi $B$ est régulier.
(On peut aussi utiliser le résultat légèrement plus général
Algèbre, lemme \ref{algebra-lemma-flat-over-regular-with-regular-fibre}.)
\end{proof}


\begin{proposition}
\label{proposition-etale-reduced}
Soient $A$, $B$ des anneaux locaux noethériens.
Soit $f : A \to B$ un homomorphisme local \'etale.
Alors $A$ est réduit si et seulement si $B$ l'est.
\end{proposition}

\begin{proof}
La fidèle platitude de $A \to B$ montre clairement que, si $B$ est réduit,
$A$ l'est aussi. Voir également Algèbre, lemme \ref{algebra-lemma-descent-reduced}.
Réciproquement, supposons $A$ réduit. Par hypothèse, $B$ est un localisé
d'une $A$-algèbre de type fini $B'$ en un idéal premier $\mathfrak q$.
Quitte à remplacer $B'$ par un localisé, nous pouvons supposer que $B'$
est \'etale sur $A$, voir le lemme \ref{lemma-characterize-etale-Noetherian}.
Alors Algèbre, lemme \ref{algebra-lemma-reduced-goes-up}, s'applique à
$A \to B'$ et $B'$ est réduit. Par conséquent, $B$ est réduit.
\end{proof}

\begin{remark}
\label{remark-technicality-needed}
Le résultat sur la ``propriété d'être réduit'' ne vaut pas avec une définition plus
faible des homomorphismes locaux \'etales $A \to B$, dans laquelle on
supprime l'hypothèse que $B$ est essentiellement de type fini sur $A$.
En effet, il peut arriver qu'un anneau local noethérien intègre $A$ ait un
complété non réduit $A^\wedge$, voir
Exemples, section \ref{examples-section-local-completion-nonreduced}.
Mais l'homomorphisme d'anneaux $A \to A^\wedge$ est plat, et $\mathfrak m_AA^\wedge$
est l'idéal maximal de $A^\wedge$ et, bien sûr, $A$ et $A^\wedge$ ont
le même corps résiduel. C'est pourquoi il importe de ne considérer
cette notion que pour des extensions d'anneaux essentiellement de type fini
(ou essentiellement de présentation finie si $A$ n'est pas noethérien).
\end{remark}

\begin{proposition}
\label{proposition-etale-normal}
\begin{reference}
\cite[Expose I, Theorem 9.5 part (i)]{SGA1}
\end{reference}
Soient $A$, $B$ des anneaux locaux noethériens.
Soit $f : A \to B$ un homomorphisme local \'etale.
Alors $A$ est un anneau intègre normal si et seulement si $B$ l'est.
\end{proposition}

\begin{proof}
Voir
Algèbre, lemme \ref{algebra-lemma-descent-normal},
pour la descente de la normalité. Réciproquement, supposons $A$ normal.
Par hypothèse, $B$ est un localisé d'une $A$-algèbre de type fini
$B'$ en un idéal premier $\mathfrak q$. Après avoir remplacé $B'$ par un localisé,
nous pouvons supposer que $B'$ est \'etale sur $A$, voir
le lemme \ref{lemma-characterize-etale-Noetherian}.
Alors
Algèbre, lemme \ref{algebra-lemma-normal-goes-up},
s'applique à $A \to B'$ et nous concluons que $B'$ est normal.
Par conséquent, $B$ est un anneau intègre normal.
\end{proof}

\noindent
Les propositions précédentes donnent une idée de la raison pour laquelle nous voulons considérer
les morphismes \'etales comme des ``isomorphismes locaux''. Une autre propriété qui indique
très clairement que nous avons la ``bonne'' définition est le fait que,
pour les $\mathbf{C}$-schémas de type fini, un morphisme est \'etale si et seulement si
le morphisme associé entre espaces analytiques (les points à valeurs dans $\mathbf{C}$ munis
de la topologie complexe) est un isomorphisme local au sens analytique (une immersion
ouverte localement sur la source). Ce fait peut être démontré à l'aide du
théorème de structure et du fait que l'analytification commute à la
formation des anneaux locaux complétés -- les détails sont laissés au lecteur.








\section{Descente des morphismes \'etales}
\label{section-descending-etale}

\noindent
Pour comprendre le langage employé dans cette section, nous invitons
le lecteur à consulter
Descente, section \ref{descent-section-descent-datum}.
Soit $f : X \to S$ un morphisme de schémas. Considérons le
foncteur de changement de base
\begin{equation}
\label{equation-descent-etale}
\text{schémas }U\text{ \'etales sur }S \longrightarrow
\begin{matrix}
\text{données de descente }(V, \varphi)\text{ relativement à }X/S \\
\text{ avec }V\text{ \'etale sur }X
\end{matrix}
\end{equation}
qui envoie $U$ sur la donnée de descente canonique $(X \times_S U, can)$.

\begin{lemma}
\label{lemma-faithful}
Si $f : X \to S$ est surjectif, alors le foncteur
(\ref{equation-descent-etale}) est fidèle.
\end{lemma}

\begin{proof}
Soient $a, b : U_1 \to U_2$ deux morphismes entre des schémas \'etales sur $S$.
Supposons que les changements de base de $a$ et $b$ par $X$ coïncident.
Nous devons montrer que $a = b$.
D'après la proposition \ref{proposition-equality}, il suffit de
montrer que $a$ et $b$ coïncident sur les points et les corps résiduels.
C'est clair, car, pour tout $u \in U_1$, nous pouvons trouver un point
$v \in X \times_S U_1$ s'envoyant sur $u$.
\end{proof}

\begin{lemma}
\label{lemma-fully-faithful}
Supposons que $f : X \to S$ soit submersif et que tout changement de base \'etale
de $f$ soit submersif. Alors le foncteur
(\ref{equation-descent-etale}) est pleinement fidèle.
\end{lemma}

\begin{proof}
D'après le lemme \ref{lemma-faithful}, le foncteur est fidèle.
Soient $U_1 \to S$ et $U_2 \to S$ des morphismes \'etales
et soit $a : X \times_S U_1 \to X \times_S U_2$ un
morphisme compatible avec les données de descente canoniques.
Nous allons montrer que $a$ est le changement de base d'un morphisme $U_1 \to U_2$.

\medskip\noindent
Soit $U'_2 \subset U_2$ un sous-schéma ouvert. Considérons
$W = a^{-1}(X \times_S U'_2)$. C'est un sous-schéma ouvert
de $X \times_S U_1$ compatible avec la donnée de descente
canonique sur $V_1 = X \times_S U_1$. Cela signifie que les
deux images réciproques de $W$ par les projections
$V_1 \times_{U_1} V_1 \to V_1$ coïncident. Puisque $V_1 \to U_1$
est surjectif (étant le changement de base de $X \to S$), nous en déduisons
que $W$ est l'image réciproque d'une partie $U'_1 \subset U_1$.
Puisque $W$ est ouvert, notre hypothèse sur $f$ implique que $U'_1 \subset U_1$
est ouvert.

\medskip\noindent
Soit $U_2 = \bigcup U_{2, i}$ un recouvrement par des ouverts affines.
D'après le résultat du paragraphe précédent, nous obtenons un recouvrement ouvert
$U_1 = \bigcup U_{1, i}$ tel que
$X \times_S U_{1, i} = a^{-1}(X \times_S U_{2, i})$.
Si nous pouvons montrer qu'il existe un morphisme $U_{1, i} \to U_{2, i}$
dont le changement de base est le morphisme
$a_i : X \times_S U_{1, i} \to X \times_S U_{2, i}$,
alors nous pouvons recoller ces morphismes en un morphisme $U_1 \to U_2$
(en utilisant la fidélité). Nous nous ramenons ainsi au cas où
$U_2$ est affine. En particulier, $U_2 \to S$ est séparé
(Schémas, lemme \ref{schemes-lemma-compose-after-separated}).

\medskip\noindent
Supposons que $U_2 \to S$ soit séparé. Alors le graphe $\Gamma_a$ de $a$
est un sous-schéma fermé de
$$
V = (X \times_S U_1) \times_X (X \times_S U_2) = X \times_S U_1 \times_S U_2
$$
d'après Schémas, lemme \ref{schemes-lemma-semi-diagonal}.
D'autre part, le graphe est ouvert, par exemple
parce que c'est une section d'un morphisme \'etale
(proposition \ref{proposition-properties-sections}).
Comme $a$ est un morphisme de données de descente, les deux images inverses de
$\Gamma_a \subset V$ par les projections
$V \times_{U_1 \times_S U_2} V \to V$ sont les mêmes.
Ainsi, en raisonnant comme au deuxième alinéa de la démonstration, nous
trouvons un sous-schéma ouvert et fermé $\Gamma \subset U_1 \times_S U_2$
dont le changement de base à $X$ donne $\Gamma_a$. Alors
$\Gamma \to U_1$ est un morphisme \'etale dont le changement de base
à $X$ est un isomorphisme. Cela signifie que $\Gamma \to U_1$
est universellement bijectif, donc est un isomorphisme
d'après le théorème \ref{theorem-etale-radicial-open}.
Ainsi, $\Gamma$ est le graphe d'un morphisme $U_1 \to U_2$
et le changement de base de ce morphisme est $a$, comme voulu.
\end{proof}

\begin{lemma}
\label{lemma-fully-faithful-cases}
Soit $f : X \to S$ un morphisme de schémas. Dans les cas suivants,
le foncteur (\ref{equation-descent-etale}) est pleinement fidèle :
\begin{enumerate}
\item $f$ est surjectif et universellement fermé
(par exemple fini, entier ou propre),
\item $f$ est surjectif et universellement ouvert
(par exemple localement de présentation finie et plat, lisse ou étale),
\item $f$ est surjectif, quasi-compact et plat.
\end{enumerate}
\end{lemma}

\begin{proof}
Cela résulte du lemme \ref{lemma-fully-faithful}.
Par exemple, une application fermée surjective d'espaces topologiques
est submersive (Topologie, lemme
\ref{topology-lemma-closed-morphism-quotient-topology}).
Les morphismes finis, entiers et propres sont universellement fermés, voir
Morphismes, lemmes \ref{morphisms-lemma-integral-universally-closed} et
\ref{morphisms-lemma-finite-proper}, et
la définition \ref{morphisms-definition-proper}.
D'autre part, une application ouverte surjective d'espaces topologiques
est submersive (Topologie, lemme
\ref{topology-lemma-open-morphism-quotient-topology}).
Les morphismes plats localement de présentation finie, lisses et \'etales sont
universellement ouverts, voir
Morphismes, lemmes \ref{morphisms-lemma-fppf-open},
\ref{morphisms-lemma-smooth-open} et
\ref{morphisms-lemma-etale-open}.
Le cas des morphismes surjectifs, quasi-compacts et plats résulte
de Morphismes, lemme \ref{morphisms-lemma-fpqc-quotient-topology}.
\end{proof}

\begin{lemma}
\label{lemma-reduce-to-affine}
Soit $f : X \to S$ un morphisme de schémas.
Soit $(V, \varphi)$ une donnée de descente relativement à $X/S$
avec $V \to X$ \'etale. Soit $S = \bigcup S_i$ un
recouvrement ouvert. Supposons que
\begin{enumerate}
\item la donnée de descente $(V, \varphi)$ obtenue par changement de base
à $X \times_S S_i/S_i$ soit effective,
\item le foncteur (\ref{equation-descent-etale})
pour $X \times_S (S_i \cap S_j) \to (S_i \cap S_j)$ soit pleinement fidèle, et
\item le foncteur (\ref{equation-descent-etale})
pour $X \times_S (S_i \cap S_j \cap S_k) \to (S_i \cap S_j \cap S_k)$
soit fidèle.
\end{enumerate}
Alors $(V, \varphi)$ est effective.
\end{lemma}

\begin{proof}
(Rappelons que les changements de base des données de descente sont définis dans
Descente, définition \ref{descent-definition-pullback-functor}.)
Posons $X_i = X \times_S S_i$. Notons $(V_i, \varphi_i)$ la donnée obtenue
par changement de base à partir de $(V, \varphi)$ sur $X_i/S_i$.
D'après l'hypothèse (1), nous pouvons trouver un morphisme \'etale $U_i \to S_i$
muni d'un isomorphisme $X_i \times_{S_i} U_i \to V_i$ compatible avec
$can$ et $\varphi_i$. D'après l'hypothèse (2), nous obtenons des isomorphismes
$\psi_{ij} : U_i \times_{S_i} (S_i \cap S_j) \to
U_j \times_{S_j} (S_i \cap S_j)$.
D'après l'hypothèse (3), ces isomorphismes satisfont à la condition de cocycle,
de sorte que $(U_i, \psi_{ij})$ est une donnée de descente pour le
recouvrement de Zariski $\{S_i \to S\}$. Alors Descente, lemme
\ref{descent-lemma-Zariski-refinement-coverings-equivalence}
(qui n'est essentiellement qu'une reformulation de
Schémas, section \ref{schemes-section-glueing-schemes})
affirme qu'il existe un morphisme de schémas $U \to S$
et des isomorphismes $U \times_S S_i \to U_i$ compatibles
avec $\psi_{ij}$. Les isomorphismes $U \times_S S_i \to U_i$
déterminent des isomorphismes correspondants $X_i \times_S U \to V_i$
qui se recollent en un morphisme $X \times_S U \to V$ compatible
avec la donnée de descente canonique et $\varphi$.
\end{proof}

\begin{lemma}
\label{lemma-split-henselian}
Soit $(A, I)$ un couple hensélien. Soit $U \to \Spec(A)$ un morphisme
quasi-compact, séparé et \'etale tel que
$U \times_{\Spec(A)} \Spec(A/I) \to \Spec(A/I)$ soit fini.
Alors
$$
U = U_{fin} \amalg U_{away}
$$
où $U_{fin} \to \Spec(A)$ est fini et $U_{away}$ ne possède
aucun point au-dessus de $Z$.
\end{lemma}

\begin{proof}
D'après le théorème principal de Zariski, le schéma $U$ est quasi-affine.
En fait, nous pouvons trouver une immersion ouverte $U \to T$, avec $T$ affine et
$T \to \Spec(A)$ fini, voir Compléments sur les morphismes, lemme
\ref{more-morphisms-lemma-quasi-finite-separated-pass-through-finite}.
Écrivons $Z = \Spec(A/I)$ et notons $U_Z \to T_Z$ le changement de base.
Comme $U_Z \to Z$ est fini, nous voyons que $U_Z \to T_Z$ est fermé
aussi bien qu'ouvert. Par conséquent, d'après
Compléments sur l'algèbre, lemme \ref{more-algebra-lemma-characterize-henselian-pair},
nous obtenons une unique décomposition $T = T' \amalg T''$ telle que $T'_Z = U_Z$.
Posons $U_{fin} = U \cap T'$ et $U_{away} = U \cap T''$. Puisque
$T'_Z \subset U_Z$, nous voyons que tous les points fermés de $T'$ sont dans $U$,
donc que $T' \subset U$, puis que $U_{fin} = T'$, et donc que $U_{fin} \to \Spec(A)$
est fini. Nous omettons la démonstration
de l'unicité de la décomposition.
\end{proof}

\begin{proposition}
\label{proposition-effective}
Soit $f : X \to S$ un morphisme entier surjectif.
Le foncteur (\ref{equation-descent-etale}) induit une équivalence
$$
\begin{matrix}
\text{schémas quasi-compacts,}\\
\text{séparés et \'etales sur }S
\end{matrix}
\longrightarrow
\begin{matrix}
\text{données de descente }(V, \varphi)\text{ relativement à }X/S\text{ avec}\\
V\text{ quasi-compact, séparé et \'etale sur }X
\end{matrix}
$$
\end{proposition}

\begin{proof}
D'après le lemme \ref{lemma-fully-faithful-cases}, le
foncteur (\ref{equation-descent-etale})
est pleinement fidèle, et cela reste vrai après tout
changement de base $S \to S'$. Soit $(V, \varphi)$ une donnée de descente
relativement à $X/S$, avec $V \to X$ quasi-compact, séparé et \'etale.
Nous pouvons utiliser le lemme \ref{lemma-reduce-to-affine}
pour voir qu'il suffit de prouver l'effectivité
localement pour la topologie de Zariski sur $S$. En particulier, nous pouvons et allons
supposer que $S$ est affine.

\medskip\noindent
Si $S$ est affine, nous pouvons trouver un ensemble ordonné filtrant $\Lambda$ et
un système projectif $X_\lambda \to S_\lambda$
de morphismes finis de schémas affines de type fini sur
$\Spec(\mathbf{Z})$ tel que $(X \to S) = \lim (X_\lambda \to S_\lambda)$.
Voir Algèbre, lemme \ref{algebra-lemma-limit-integral}.
Puisque les limites commutent entre elles, nous en déduisons que
$X \times_S X = \lim X_\lambda \times_{S_\lambda} X_\lambda$
et que
$X \times_S X \times_S X = \lim
X_\lambda \times_{S_\lambda} X_\lambda \times_{S_\lambda} X_\lambda$.
Remarquons que $V \to X$ est un morphisme de présentation finie.
En utilisant Limites, lemme \ref{limits-lemma-descend-finite-presentation},
nous pouvons trouver un $\lambda$ et une donnée de descente $(V_\lambda, \varphi_\lambda)$
relativement à $X_\lambda/S_\lambda$ dont le changement de base à $X/S$ est
$(V, \varphi)$. Il suffit bien sûr de montrer que
$(V_\lambda, \varphi_\lambda)$ est effective. Notons que $V_\lambda$
est quasi-compact par construction.
Quitte à augmenter $\lambda$, nous pouvons supposer
que $V_\lambda \to X_\lambda$ est séparé et \'etale, voir
Limites, lemme \ref{limits-lemma-descend-separated-finite-presentation} et
\ref{limits-lemma-descend-etale}.
Ainsi, nous pouvons supposer que $f$ est fini surjectif et que
$S$ est affine de type fini sur $\mathbf{Z}$.

\medskip\noindent
Considérons un ouvert $S' \subset S$ tel que la donnée $(V', \varphi')$ obtenue
par changement de base à partir de $(V, \varphi)$ sur $X' = X \times_S S'$ soit effective. Nous prouverons ci-dessous
que $S' \not = S$ implique l'existence d'un ouvert strictement plus grand sur
lequel la donnée de descente est effective. Comme $S$ est noethérien (et que son
espace topologique sous-jacent est donc noethérien), cela achèvera la démonstration.
Soit $\xi \in S$ un point générique d'une composante irréductible du
fermé $Z = S \setminus S'$.
Si $\xi \in S'' \subset S$ est un ouvert sur lequel la donnée de descente est
effective, alors celle-ci est effective sur
$S' \cup S''$ par l'argument de recollement du premier alinéa. Ainsi,
dans la suite de la démonstration, nous pouvons remplacer $S$ par un voisinage ouvert affine
de $\xi$.

\medskip\noindent
Après un premier remplacement de ce type, nous pouvons supposer que $Z$ est irréductible
de point générique $Z$. Munissons $Z$ de la structure réduite induite
de sous-schéma fermé. Après une nouvelle restriction, nous pouvons supposer que
$X_Z = X \times_S Z = f^{-1}(Z) \to Z$ est plat, voir
Morphismes, proposition \ref{morphisms-proposition-generic-flatness}.
Soit $(V_Z, \varphi_Z)$ la donnée de descente obtenue par changement de base à $X_Z/Z$.
D'après Compléments sur les morphismes, lemme
\ref{more-morphisms-lemma-separated-locally-quasi-finite-morphisms-fppf-descend},
cette donnée de descente est effective, et nous obtenons un morphisme \'etale
$U_Z \to Z$ dont le changement de base est isomorphe à $V_Z$ d'une manière
compatible avec les données de descente.
Bien sûr, $U_Z \to Z$ est quasi-compact et séparé
(Descente, lemmes \ref{descent-lemma-descending-property-quasi-compact} et
\ref{descent-lemma-descending-property-separated}).
Ainsi, après une nouvelle restriction, nous pouvons supposer
que $U_Z \to Z$ est fini, voir
Morphismes, lemme \ref{morphisms-lemma-generically-finite}.

{\emergencystretch=.5em\medskip\noindent
Soit $S = \Spec(A)$ et soit $I \subset A$ l'idéal premier correspondant
à $Z \subset S$. Soit $(A^h, IA^h)$ l'hensélisation du couple
$(A, I)$.\linebreak Notons $S^h = \Spec(A^h)$ et $Z^h = V(IA^h) \cong Z$.
Nous affirmons qu'il suffit de montrer l'effectivité après changement de base à
$S^h$. En effet, $\{S^h \to S, S' \to S\}$ est un recouvrement fpqc
($A \to A^h$ est plat d'après Compléments sur l'algèbre, lemme
\ref{more-algebra-lemma-henselization-flat}) et,
d'après Compléments sur les morphismes, lemme
\ref{more-morphisms-lemma-separated-locally-quasi-finite-morphisms-fppf-descend},
nous disposons de la descente fpqc pour les morphismes \'etales séparés.
En effet, si $U^h \to S^h$ et $U' \to S'$ sont les objets
correspondant aux images inverses $(V^h, \varphi^h)$ et
$(V', \varphi')$, alors les isomorphismes requis\par}
$$
U^h \times_S S^h \to S^h \times_S V^h
\quad\text{et}\quad
U^h \times_S S' \to S^h \times_S U'
$$
sont obtenus grâce à la pleine fidélité signalée au premier
alinéa. Nous nous ramenons ainsi à la situation décrite dans
l'alinéa suivant.

\medskip\noindent
Ici, $S = \Spec(A)$, $Z = V(I)$, $S' = S \setminus Z$, où
$(A, I)$ est un couple hensélien ; nous avons $U' \to S'$ correspondant
à la donnée de descente $(V', \varphi')$ et un morphisme \'etale fini
$U_Z \to Z$ correspondant à la donnée de descente
$(V_Z, \varphi_Z)$. Nous ne disposons plus de l'hypothèse que $A$ est de type fini
sur $\mathbf{Z}$, mais le reste de l'argument n'utilisera même pas
que $A$ est noethérien.
D'après Compléments sur l'algèbre, lemme \ref{more-algebra-lemma-finite-etale-equivalence},
nous pouvons trouver un morphisme \'etale fini $U_{fin} \to S$ dont la
restriction à $Z$ est isomorphe à $U_Z \to Z$.
Écrivons $X = \Spec(B)$ et $Y = V(IB)$. Comme $(B, IB)$ est un couple hensélien
(Compléments sur l'algèbre, lemme \ref{more-algebra-lemma-integral-over-henselian-pair})
et comme la restriction de $V \to X$ à $Y$
est finie (car elle est le changement de base de $U_Z \to Z$), nous voyons
qu'il existe une décomposition canonique en union disjointe
$$
V = V_{fin} \amalg V_{away}
$$
où $V_{fin} \to X$ est fini et où $V_{away}$ ne possède aucun
point au-dessus de $Y$. Voir le lemme \ref{lemma-split-henselian}.
En utilisant l'unicité de cette décomposition sur $X \times_S X$,
nous voyons que $\varphi$ la respecte, et nous obtenons
$$
(V, \varphi) = (V_{fin}, \varphi_{fin}) \amalg (V_{away}, \varphi_{away})
$$
dans la catégorie des données de descente.
D'après Compléments sur l'algèbre, lemme \ref{more-algebra-lemma-finite-etale-equivalence},
il existe un unique isomorphisme
$$
X \times_S U_{fin} \longrightarrow V_{fin}
$$
compatible avec l'isomorphisme donné $Y \times_Z U_Z \to V \times_X Y$
sur $Y$.
Par unicité, nous voyons que cet isomorphisme est compatible
avec les données de descente, c'est-à-dire que
$(X \times_S U_{fin}, can) \cong (V_{fin}, \varphi_{fin})$.
Notons $U'_{fin} = U_{fin} \times_S S'$. Par pleine fidélité,
nous obtenons un morphisme $U'_{fin} \to U'$ qui est
l'inclusion d'un sous-schéma ouvert (et fermé).
Posons alors $U = U_{fin} \amalg_{U'_{fin}} U'$ (recollement de schémas comme
dans Schémas, section \ref{schemes-section-glueing-schemes}).
Les morphismes $X \times_S U_{fin} \to V$ et
$X \times_S U' \to V$ se recollent en un morphisme $X \times_S U \to V$
qui est l'isomorphisme recherché.
\end{proof}





\section{Diviseurs à croisements normaux}
\label{section-normal-crossings}

\noindent
Voici la définition.

\begin{definition}
\label{definition-strict-normal-crossings}
Soit $X$ un schéma localement noethérien. Un
{\it diviseur à croisements strictement normaux}
sur $X$ est un diviseur de Cartier effectif $D \subset X$ tel que,
pour tout $p \in D$, l'anneau local $\mathcal{O}_{X, p}$ soit régulier
et qu'il existe un système régulier de paramètres
$x_1, \ldots, x_d \in \mathfrak m_p$ et $1 \leq r \leq d$
tel que $D$ soit défini par $x_1 \ldots x_r$ dans $\mathcal{O}_{X, p}$.
\end{definition}

\noindent
Nous rencontrons souvent des diviseurs de Cartier effectifs $E$ sur des schémas
localement noethériens $X$ tels qu'il existe un diviseur à croisements strictement normaux $D$
avec $E \subset D$ ensemblistement.
Dans ce cas, nous avons
$E = \sum a_i D_i$, avec $a_i \geq 0$, où $D = \bigcup_{i \in I} D_i$
est la décomposition de $D$ en ses composantes irréductibles.
Remarquons que $D' = \bigcup_{a_i > 0} D_i$ est un diviseur à croisements
normaux stricts tel que $E = D'$ ensemblistement.
Lorsque ce qui précède se produit, nous dirons que
$E$ est {\it à support contenu dans un diviseur à croisements strictement normaux}.

\begin{lemma}
\label{lemma-strict-normal-crossings}
Soit $X$ un schéma localement noethérien. Soit $D \subset X$ un
diviseur de Cartier effectif. Soient $D_i \subset D$, $i \in I$, ses
composantes irréductibles considérées comme des sous-schémas fermés réduits de $X$.
Les conditions suivantes sont équivalentes :
\begin{enumerate}
\item $D$ est un diviseur à croisements strictement normaux, et
\item $D$ est réduit, chaque $D_i$ est un diviseur de Cartier effectif, et,
pour $J \subset I$ fini, l'intersection schématique
$D_J = \bigcap_{j \in J} D_j$ est un schéma
régulier dont chaque composante irréductible est de
codimension $|J|$ dans $X$.
\end{enumerate}
\end{lemma}

\begin{proof}
Supposons que $D$ soit un diviseur à croisements strictement normaux. Choisissons $p \in D$
et un système régulier de paramètres $x_1, \ldots, x_d \in \mathfrak m_p$
ainsi que $1 \leq r \leq d$ comme dans la
définition \ref{definition-strict-normal-crossings}.
Comme $\mathcal{O}_{X, p}/(x_i)$ est un anneau local régulier
(et en particulier un anneau intègre), nous voyons que les composantes irréductibles
$D_1, \ldots, D_r$ de $D$ passant par $p$ correspondent $1$ à $1$
aux idéaux premiers de hauteur un $(x_1), \ldots, (x_r)$ de $\mathcal{O}_{X, p}$.
D'après Algèbre, lemme \ref{algebra-lemma-regular-ring-CM},
nous trouvons que l'intersection $D_{i_1} \cap \ldots \cap D_{i_s}$
est de codimension $s$ dans un voisinage ouvert de $p$
et que son anneau local en $p$ est régulier.
Comme cela vaut pour tout $p \in D$, nous concluons que (2) est satisfaite.

\medskip\noindent
Supposons (2). Soit $p \in D$. Comme $\mathcal{O}_{X, p}$ est de dimension
finie, nous voyons que $p$ peut appartenir à au plus
$\dim(\mathcal{O}_{X, p})$ des composantes $D_i$.
Disons que $p \in D_1, \ldots, D_r$ pour un certain $r \geq 1$.
Soient $x_1, \ldots, x_r \in \mathfrak m_p$ des équations locales
de $D_1, \ldots, D_r$. Alors $x_1$ est un non diviseur de zéro dans $\mathcal{O}_{X, p}$
et $\mathcal{O}_{X, p}/(x_1) = \mathcal{O}_{D_1, p}$ est régulier.
Ainsi, $\mathcal{O}_{X, p}$ est régulier, voir
Algèbre, lemme \ref{algebra-lemma-regular-mod-x}.
Comme $D_1 \cap \ldots \cap D_r$ est un schéma régulier (donc normal),
c'est une union disjointe de ses composantes irréductibles
(Propriétés, lemme \ref{properties-lemma-normal-Noetherian}).
Soit $Z \subset D_1 \cap \ldots \cap D_r$
la composante irréductible contenant $p$.
Alors $\mathcal{O}_{Z, p} = \mathcal{O}_{X, p}/(x_1, \ldots, x_r)$
est régulier de codimension $r$ (notons que, puisque nous savons déjà
que $\mathcal{O}_{X, p}$ est régulier et donc de Cohen-Macaulay,
il n'y a aucune ambiguïté sur la codimension, car l'anneau est caténaire, voir
Algèbre, lemmes \ref{algebra-lemma-regular-ring-CM} et
\ref{algebra-lemma-CM-dim-formula}).
Ainsi, $\dim(\mathcal{O}_{Z, p}) = \dim(\mathcal{O}_{X, p}) - r$.
Choisissons des éléments supplémentaires $x_{r + 1}, \ldots, x_n \in \mathfrak m_p$
dont les images forment un système minimal de générateurs de $\mathfrak m_{Z, p}$.
Alors $\mathfrak m_p = (x_1, \ldots, x_n)$ d'après le lemme de Nakayama,
et nous voyons que $D$ est un diviseur à croisements normaux.
\end{proof}

\begin{lemma}
\label{lemma-smooth-pullback-strict-normal-crossings}
\begin{slogan}
Le tiré en arrière d'un diviseur à croisements strictement normaux par un morphisme
lisse est un diviseur à croisements strictement normaux.
\end{slogan}
Soit $X$ un schéma localement noethérien. Soit $D \subset X$ un
diviseur à croisements strictement normaux. Si $f : Y \to X$ est un morphisme
lisse de schémas, alors le tiré en arrière $f^*D$ est un
diviseur à croisements strictement normaux sur $Y$.
\end{lemma}

\begin{proof}
Comme $f$ est plat, le tiré en arrière est défini par
Diviseurs, lemme \ref{divisors-lemma-pullback-effective-Cartier-defined},
donc l'énoncé a un sens.
Soit $q \in f^*D$ d'image $p \in D$. Choisissons un système régulier
de paramètres $x_1, \ldots, x_d \in \mathfrak m_p$
et $1 \leq r \leq d$ comme dans la
définition \ref{definition-strict-normal-crossings}.
Comme $f$ est lisse, l'homomorphisme d'anneaux locaux
$\mathcal{O}_{X, p} \to \mathcal{O}_{Y, q}$ est plat
et l'anneau de la fibre
$$
\mathcal{O}_{Y, q}/\mathfrak m_p \mathcal{O}_{Y, q} =
\mathcal{O}_{Y_p, q}
$$
est un anneau local régulier (voir par exemple
Algèbre, lemme \ref{algebra-lemma-characterize-smooth-over-field}).
Choisissons $y_1, \ldots, y_n \in \mathfrak m_q$ dont les images forment un système
régulier de paramètres dans $\mathcal{O}_{Y_p, q}$.
Alors $x_1, \ldots, x_d, y_1, \ldots, y_n$ engendrent l'idéal
maximal $\mathfrak m_q$. Ainsi, $\mathcal{O}_{Y, q}$
est un anneau local régulier de dimension
$d + n$ d'après Algèbre, lemme \ref{algebra-lemma-dimension-base-fibre-equals-total},
et $x_1, \ldots, x_d, y_1, \ldots, y_n$
est un système régulier de paramètres. Comme $f^*D$ est défini
par $x_1 \ldots x_r$ dans $\mathcal{O}_{Y, q}$, nous concluons
que le lemme est démontré.
\end{proof}

\noindent
Voici la définition d'un diviseur à croisements normaux.

\begin{definition}
\label{definition-normal-crossings}
Soit $X$ un schéma localement noethérien. Un {\it diviseur à croisements normaux}
sur $X$ est un diviseur de Cartier effectif $D \subset X$ tel que, pour
tout $p \in D$, il existe un morphisme \'etale $U \to X$ dont l'image
contient $p$ et tel que $D \times_X U$ soit un
diviseur à croisements strictement normaux sur $U$.
\end{definition}

\noindent
Par exemple, $D = V(x^2 + y^2)$ est un diviseur à croisements normaux
(mais non strict) sur
$\Spec(\mathbf{R}[x, y])$, car, après changement de base au
recouvrement \'etale $\Spec(\mathbf{C}[x, y])$, nous obtenons $(x - iy)(x + iy) = 0$.

\begin{lemma}
\label{lemma-smooth-pullback-normal-crossings}
\begin{slogan}
Le tiré en arrière d'un diviseur à croisements normaux par un morphisme
lisse est un diviseur à croisements normaux.
\end{slogan}
Soit $X$ un schéma localement noethérien. Soit $D \subset X$ un
diviseur à croisements normaux. Si $f : Y \to X$ est un morphisme
lisse de schémas, alors le tiré en arrière $f^*D$ est un
diviseur à croisements normaux sur $Y$.
\end{lemma}

\begin{proof}
Comme $f$ est plat, le tiré en arrière est défini par
Diviseurs, lemme \ref{divisors-lemma-pullback-effective-Cartier-defined},
donc l'énoncé a un sens.
Soit $q \in f^*D$ d'image $p \in D$.
Choisissons un morphisme \'etale $U \to X$ dont l'image contient $p$
tel que $D \times_X U \subset U$ soit un diviseur à croisements normaux
stricts, comme dans la définition \ref{definition-normal-crossings}.
Posons $V = Y \times_X U$. Alors $V \to Y$ est \'etale comme changement
de base de $U \to X$
(Morphismes, lemme \ref{morphisms-lemma-base-change-etale}),
et le tiré en arrière $D \times_X V$ est un diviseur à croisements normaux
stricts sur $V$ d'après le lemme \ref{lemma-smooth-pullback-strict-normal-crossings}.
Nous avons donc vérifié la condition de la
définition \ref{definition-normal-crossings}
pour $q \in f^*D$, ce qui permet de conclure.
\end{proof}

\begin{lemma}
\label{lemma-characterize-normal-crossings-normalization}
Soit $X$ un schéma localement noethérien. Soit $D \subset X$ un sous-schéma
fermé. Les conditions suivantes sont équivalentes :
\begin{enumerate}
\item $D$ est un diviseur à croisements normaux dans $X$,
\item $D$ est réduit, la normalisation $\nu : D^\nu \to D$ est non ramifiée,
et, pour tout $n \geq 1$, le schéma
$$
Z_n = D^\nu \times_D \ldots \times_D D^\nu
\setminus \{(p_1, \ldots, p_n) \mid p_i = p_j\text{ pour certains }i\not = j\}
$$
est régulier ; le morphisme $Z_n \to X$ est d'intersection complète locale
et son faisceau conormal est localement libre de rang $n$.
\end{enumerate}
\end{lemma}

\begin{proof}
Expliquons d'abord comment interpréter la condition (2).
La diagonale d'un morphisme non ramifié est ouverte
(Morphismes, lemme \ref{morphisms-lemma-diagonal-unramified-morphism}).
D'autre part, $D^\nu \to D$ est séparé, donc la
diagonale $D^\nu \to D^\nu \times_D D^\nu$ est fermée.
Ainsi, $Z_n$ est un sous-schéma ouvert et fermé de
$D^\nu \times_D \ldots \times_D D^\nu$. D'autre part,
$Z_n \to X$ est non ramifié, car c'est la composée
$$
Z_n \to D^\nu \times_D \ldots \times_D D^\nu \to \ldots \to
D^\nu \times_D D^\nu \to D^\nu \to D \to X
$$
et chacune des flèches est non ramifiée.
Comme un morphisme non ramifié est formellement non ramifié
(Compléments sur les morphismes, lemme
\ref{more-morphisms-lemma-unramified-formally-unramified}),
nous avons un faisceau conormal
$\mathcal{C}_n = \mathcal{C}_{Z_n/X}$ du morphisme $Z_n \to X$, voir
Compléments sur les morphismes, définition
\ref{more-morphisms-definition-universal-thickening}.

\medskip\noindent
La formation de la normalisation commute à la localisation \'etale d'après
Compléments sur les morphismes, lemme \ref{more-morphisms-lemma-normalization-and-smooth}.
La régularité des anneaux locaux, le fait
qu'un morphisme soit non ramifié ou localement
d'intersection complète, ou le fait qu'un morphisme soit
non ramifié et possède un faisceau conormal qui soit
localement libre d'un rang donné, se vérifient localement pour la topologie \'etale (voir
Compléments sur l'algèbre, lemme \ref{more-algebra-lemma-regular-etale-extension},
Descente, lemme \ref{descent-lemma-descending-property-unramified},
Compléments sur les morphismes, lemme \ref{more-morphisms-lemma-descending-property-lci}
et
Descente, lemme \ref{descent-lemma-finite-locally-free-descends}).

\medskip\noindent
D'après la remarque de l'alinéa précédent et la définition
des diviseurs à croisements normaux, il suffit de prouver qu'un
diviseur à croisements strictement normaux $D = \bigcup_{i \in I} D_i$
satisfait à (2). Dans ce cas, $D^\nu = \coprod D_i$
et $D^\nu \to D$ est non ramifié (la propriété d'être non ramifié
est locale sur la source, et $D_i \to D$ est une immersion
fermée qui est non ramifiée). De même, $Z_1 = D^\nu \to X$
est un morphisme d'intersection complète locale, car nous pouvons
le vérifier localement sur la source et chaque morphisme $D_i \to X$
est une immersion régulière, puisque c'est l'inclusion d'un diviseur de Cartier
(voir le lemme \ref{lemma-strict-normal-crossings} et
Compléments sur les morphismes, lemme \ref{more-morphisms-lemma-regular-immersion-lci}).
Comme un diviseur de Cartier effectif possède un faisceau
conormal inversible, nous concluons que la condition portant sur le
faisceau conormal est satisfaite.
De même, le schéma $Z_n$, pour $n \geq 2$, est l'union disjointe
des schémas $D_J = \bigcap_{j \in J} D_j$, où $J \subset I$
parcourt les sous-ensembles de cardinal $n$. Comme $D_J \to X$ est
une immersion régulière de codimension $n$
(par la définition des croisements normaux stricts et le
fait que nous pouvons le vérifier sur les germes d'après
Diviseurs, lemme \ref{divisors-lemma-Noetherian-scheme-regular-ideal}),
il s'ensuit de la même manière que $Z_n \to X$ possède les propriétés
requises. Nous omettons quelques détails.

\medskip\noindent
Supposons (2). Soit $p \in D$. Comme $D^\nu \to D$ est non ramifié, il est
fini (d'après Morphismes, lemme \ref{morphisms-lemma-finite-integral}).
Ainsi, $D^\nu \to X$ est fini non ramifié.
D'après le lemme \ref{lemma-finite-unramified-etale-local}
et par localisation \'etale (permise par la discussion
du deuxième alinéa et par la définition des diviseurs à
croisements normaux), nous nous ramenons au cas où
$D^\nu = \coprod_{i \in I} D_i$,
avec $I$ fini et $D_i \to U$ une immersion fermée.
Quitte à restreindre $X$, nous pouvons supposer
que $p \in D_i$ pour tout $i \in I$. La condition selon laquelle $Z_1 =  D^\nu \to X$ est un
morphisme non ramifié d'intersection complète locale
dont le faisceau conormal est localement libre de rang $1$
implique que $D_i \subset X$ est un diviseur de Cartier effectif, voir
Compléments sur les morphismes, lemme \ref{more-morphisms-lemma-lci} et
Diviseurs, lemme \ref{divisors-lemma-regular-immersion-noetherian}.
Pour achever la démonstration, nous pouvons supposer que $X = \Spec(A)$ est affine
et que $D_i = V(f_i)$, avec $f_i \in A$ un non diviseur de zéro.
Si $I = \{1, \ldots, r\}$, alors $p \in Z_r = V(f_1, \ldots, f_r)$.
La même référence que ci-dessus implique que
$(f_1, \ldots, f_r)$ est un idéal Koszul-régulier dans $A$.
Comme le faisceau conormal est de rang $r$, nous voyons que
$f_1, \ldots, f_r$ est un système minimal de générateurs de
l'idéal définissant $Z_r$ dans $\mathcal{O}_{X, p}$.
Cela implique que $f_1, \ldots, f_r$ est une suite régulière
dans $\mathcal{O}_{X, p}$ telle que $\mathcal{O}_{X, p}/(f_1, \ldots, f_r)$
soit régulier. Nous concluons ainsi du
lemme \ref{algebra-lemma-regular-mod-x} d'Algèbre
que la suite $f_1, \ldots, f_r$ peut être complétée en un système régulier de paramètres
dans $\mathcal{O}_{X, p}$, ce qui achève la démonstration.
\end{proof}

\begin{lemma}
\label{lemma-characterize-normal-crossings}
Soit $X$ un schéma localement noethérien. Soit $D \subset X$ un sous-schéma
fermé. Si $X$ est J-2 ou de Nagata, les conditions suivantes sont équivalentes :
\begin{enumerate}
\item $D$ est un diviseur à croisements normaux dans $X$,
\item pour tout $p \in D$, le tiré en arrière de $D$ sur le spectre de
l'hensélisé strict $\mathcal{O}_{X, p}^{sh}$
est un diviseur à croisements strictement normaux.
\end{enumerate}
\end{lemma}

\begin{proof}
L'implication (1) $\Rightarrow$ (2) est immédiate et
ne nécessite pas l'hypothèse que $X$ est J-2 ou de Nagata.
En effet, soit $p \in D$ et choisissons un voisinage \'etale
$(U, u) \to (X, p)$ tel que le tiré en arrière de $D$ soit
un diviseur à croisements strictement normaux sur $U$.
Alors $\mathcal{O}_{X, p}^{sh} = \mathcal{O}_{U, u}^{sh}$,
et nous voyons que la trace de $D$ sur $\Spec(\mathcal{O}_{U, u}^{sh})$
est définie par une partie d'un système régulier de paramètres,
puisque c'est déjà le cas dans $\mathcal{O}_{U, u}$.

\medskip\noindent
Pour prouver l'implication dans l'autre sens,
nous utiliserons le critère du
lemme \ref{lemma-characterize-normal-crossings-normalization}.
Remarquons que la formation de la normalisation $D^\nu \to D$
commute à l'hensélisation stricte, voir
Compléments sur les morphismes, lemme
\ref{more-morphisms-lemma-normalization-and-henselization}.
Si nous pouvons montrer que $D^\nu \to D$ est fini,
alors nous voyons que $D^\nu \to D$ et les schémas
$Z_n$ satisfont à toutes les propriétés voulues, car celles-ci
peuvent toutes se vérifier au niveau des anneaux locaux
(mais la finitude du morphisme $D^\nu \to D$
ne peut pas se vérifier sur les anneaux locaux).
Nous omettons les vérifications détaillées.

\medskip\noindent
Si $X$ est de Nagata, alors $D^\nu \to D$ est fini d'après
Morphismes, lemme \ref{morphisms-lemma-nagata-normalization}.

\medskip\noindent
Supposons que $X$ soit J-2. Choisissons un point $p \in D$. Nous montrerons
que $D^\nu \to D$ est fini au-dessus d'un voisinage de $p$.
Par hypothèse, il existe un système régulier de
paramètres $f_1, \ldots, f_d$ de $\mathcal{O}_{X, p}^{sh}$
et $1 \leq r \leq d$ tels que la trace de $D$ sur
$\Spec(\mathcal{O}_{X, p}^{sh})$ soit définie par $f_1 \ldots f_r$.
Alors
$$
D^\nu \times_X \Spec(\mathcal{O}_{X, p}^{sh}) = 
\coprod\nolimits_{i = 1, \ldots, r} V(f_i)
$$
Choisissons un voisinage \'etale affine
$(U, u) \to (X, p)$ tel que $f_i$ provienne d'un élément
$f_i \in \mathcal{O}_U(U)$. Posons $D_i = V(f_i) \subset U$.
L'hensélisé strict de $\mathcal{O}_{D_i, u}$
est $\mathcal{O}_{X, p}^{sh}/(f_i)$, qui est régulier.
Ainsi, $\mathcal{O}_{D_i, u}$ est régulier (par exemple d'après
Compléments sur l'algèbre, lemme \ref{more-algebra-lemma-henselization-regular}).
Comme $X$ est J-2, le lieu régulier est ouvert dans $D_i$.
Ainsi, quitte à remplacer $U$ par un ouvert de Zariski, nous pouvons supposer
que $D_i$ est régulier pour tout $i$. Il s'ensuit que
$$
\coprod\nolimits_{i = 1, \ldots, r} D_i = D^\nu \times_X U
\longrightarrow D \times_X U
$$
est le morphisme de normalisation, et il est clairement fini.
Autrement dit, nous avons trouvé
un voisinage \'etale $(U, u)$ de $(X, p)$ tel que
le changement de base de $D^\nu \to D$ à ce voisinage soit fini.
Cela implique que $D^\nu \to D$ est fini par descente
(Descente, lemme \ref{descent-lemma-descending-property-finite}),
et la démonstration est achevée.
\end{proof}






\begin{multicols}{2}[\section{Autres chapitres}]
\noindent
Préliminaires
\begin{enumerate}
\item \hyperref[introduction-section-phantom]{Introduction}
\item \hyperref[conventions-section-phantom]{Conventions}
\item \hyperref[sets-section-phantom]{Théorie des ensembles}
\item \hyperref[categories-section-phantom]{Catégories}
\item \hyperref[topology-section-phantom]{Topologie}
\item \hyperref[sheaves-section-phantom]{Faisceaux sur les espaces}
\item \hyperref[sites-section-phantom]{Sites et faisceaux}
\item \hyperref[stacks-section-phantom]{Champs}
\item \hyperref[fields-section-phantom]{Corps}
\item \hyperref[algebra-section-phantom]{Algèbre commutative}
\item \hyperref[brauer-section-phantom]{Groupes de Brauer}
\item \hyperref[homology-section-phantom]{Algèbre homologique}
\item \hyperref[derived-section-phantom]{Catégories dérivées}
\item \hyperref[simplicial-section-phantom]{Méthodes simpliciales}
\item \hyperref[more-algebra-section-phantom]{Compléments d'algèbre}
\item \hyperref[smoothing-section-phantom]{Lissage des morphismes d'anneaux}
\item \hyperref[modules-section-phantom]{Faisceaux de modules}
\item \hyperref[sites-modules-section-phantom]{Modules sur les sites}
\item \hyperref[injectives-section-phantom]{Objets injectifs}
\item \hyperref[cohomology-section-phantom]{Cohomologie des faisceaux}
\item \hyperref[sites-cohomology-section-phantom]{Cohomologie sur les sites}
\item \hyperref[dga-section-phantom]{Algèbre différentielle graduée}
\item \hyperref[dpa-section-phantom]{Algèbre à puissances divisées}
\item \hyperref[sdga-section-phantom]{Faisceaux différentiels gradués}
\item \hyperref[hypercovering-section-phantom]{Hyperrecouvrements}
\end{enumerate}
Schémas
\begin{enumerate}
\setcounter{enumi}{25}
\item \hyperref[schemes-section-phantom]{Schémas}
\item \hyperref[constructions-section-phantom]{Constructions de schémas}
\item \hyperref[properties-section-phantom]{Propriétés des schémas}
\item \hyperref[morphisms-section-phantom]{Morphismes de schémas}
\item \hyperref[coherent-section-phantom]{Cohomologie des schémas}
\item \hyperref[divisors-section-phantom]{Diviseurs}
\item \hyperref[limits-section-phantom]{Limites de schémas}
\item \hyperref[varieties-section-phantom]{Variétés}
\item \hyperref[topologies-section-phantom]{Topologies sur les schémas}
\item \hyperref[descent-section-phantom]{Descente}
\item \hyperref[perfect-section-phantom]{Catégories dérivées de schémas}
\item \hyperref[more-morphisms-section-phantom]{Compléments sur les morphismes}
\item \hyperref[flat-section-phantom]{Compléments sur la platitude}
\item \hyperref[groupoids-section-phantom]{Schémas en groupoïdes}
\item \hyperref[more-groupoids-section-phantom]{Compléments sur les schémas en groupoïdes}
\item \hyperref[etale-section-phantom]{Morphismes étales de schémas}
\end{enumerate}
Thèmes de théorie des schémas
\begin{enumerate}
\setcounter{enumi}{41}
\item \hyperref[chow-section-phantom]{Homologie de Chow}
\item \hyperref[intersection-section-phantom]{Théorie de l'intersection}
\item \hyperref[pic-section-phantom]{Schémas de Picard des courbes}
\item \hyperref[weil-section-phantom]{Théories cohomologiques de Weil}
\item \hyperref[adequate-section-phantom]{Modules adéquats}
\item \hyperref[dualizing-section-phantom]{Complexes dualisants}
\item \hyperref[duality-section-phantom]{Dualité pour les schémas}
\item \hyperref[discriminant-section-phantom]{Discriminants et différentes}
\item \hyperref[derham-section-phantom]{Cohomologie de de Rham}
\item \hyperref[local-cohomology-section-phantom]{Cohomologie locale}
\item \hyperref[algebraization-section-phantom]{Géométrie algébrique et formelle}
\item \hyperref[curves-section-phantom]{Courbes algébriques}
\item \hyperref[resolve-section-phantom]{Résolution des surfaces}
\item \hyperref[models-section-phantom]{Réduction semi-stable}
\item \hyperref[functors-section-phantom]{Foncteurs et morphismes}
\item \hyperref[equiv-section-phantom]{Catégories dérivées de variétés}
\item \hyperref[pione-section-phantom]{Groupes fondamentaux de schémas}
\item \hyperref[etale-cohomology-section-phantom]{Cohomologie étale}
\item \hyperref[crystalline-section-phantom]{Cohomologie cristalline}
\item \hyperref[proetale-section-phantom]{Cohomologie pro-étale}
\item \hyperref[relative-cycles-section-phantom]{Cycles relatifs}
\item \hyperref[more-etale-section-phantom]{Compléments de cohomologie étale}
\item \hyperref[trace-section-phantom]{La formule des traces}
\end{enumerate}
Espaces algébriques
\begin{enumerate}
\setcounter{enumi}{64}
\item \hyperref[spaces-section-phantom]{Espaces algébriques}
\item \hyperref[spaces-properties-section-phantom]{Propriétés des espaces algébriques}
\item \hyperref[spaces-morphisms-section-phantom]{Morphismes d'espaces algébriques}
\item \hyperref[decent-spaces-section-phantom]{Espaces algébriques décents}
\item \hyperref[spaces-cohomology-section-phantom]{Cohomologie des espaces algébriques}
\item \hyperref[spaces-limits-section-phantom]{Limites d'espaces algébriques}
\item \hyperref[spaces-divisors-section-phantom]{Diviseurs sur les espaces algébriques}
\item \hyperref[spaces-over-fields-section-phantom]{Espaces algébriques sur des corps}
\item \hyperref[spaces-topologies-section-phantom]{Topologies sur les espaces algébriques}
\item \hyperref[spaces-descent-section-phantom]{Descente et espaces algébriques}
\item \hyperref[spaces-perfect-section-phantom]{Catégories dérivées d'espaces}
\item \hyperref[spaces-more-morphisms-section-phantom]{Compléments sur les morphismes d'espaces}
\item \hyperref[spaces-flat-section-phantom]{Platitude sur les espaces algébriques}
\item \hyperref[spaces-groupoids-section-phantom]{Groupoïdes dans les espaces algébriques}
\item \hyperref[spaces-more-groupoids-section-phantom]{Compléments sur les groupoïdes dans les espaces}
\item \hyperref[bootstrap-section-phantom]{Amorçage}
\item \hyperref[spaces-pushouts-section-phantom]{Sommes amalgamées d'espaces algébriques}
\end{enumerate}
Thèmes de géométrie
\begin{enumerate}
\setcounter{enumi}{81}
\item \hyperref[spaces-chow-section-phantom]{Groupes de Chow des espaces}
\item \hyperref[groupoids-quotients-section-phantom]{Quotients de groupoïdes}
\item \hyperref[spaces-more-cohomology-section-phantom]{Compléments sur la cohomologie des espaces}
\item \hyperref[spaces-simplicial-section-phantom]{Espaces simpliciaux}
\item \hyperref[spaces-duality-section-phantom]{Dualité pour les espaces}
\item \hyperref[formal-spaces-section-phantom]{Espaces algébriques formels}
\item \hyperref[restricted-section-phantom]{Algébrisation des espaces formels}
\item \hyperref[spaces-resolve-section-phantom]{Résolution des surfaces revisitée}
\end{enumerate}
Théorie des déformations
\begin{enumerate}
\setcounter{enumi}{89}
\item \hyperref[formal-defos-section-phantom]{Théorie formelle des déformations}
\item \hyperref[defos-section-phantom]{Théorie des déformations}
\item \hyperref[cotangent-section-phantom]{Le complexe cotangent}
\item \hyperref[examples-defos-section-phantom]{Problèmes de déformation}
\end{enumerate}
Champs algébriques
\begin{enumerate}
\setcounter{enumi}{93}
\item \hyperref[algebraic-section-phantom]{Champs algébriques}
\item \hyperref[examples-stacks-section-phantom]{Exemples de champs}
\item \hyperref[stacks-sheaves-section-phantom]{Faisceaux sur les champs algébriques}
\item \hyperref[criteria-section-phantom]{Critères de représentabilité}
\item \hyperref[artin-section-phantom]{Axiomes d'Artin}
\item \hyperref[quot-section-phantom]{Espaces de Quot et de Hilbert}
\item \hyperref[stacks-properties-section-phantom]{Propriétés des champs algébriques}
\item \hyperref[stacks-morphisms-section-phantom]{Morphismes de champs algébriques}
\item \hyperref[stacks-limits-section-phantom]{Limites de champs algébriques}
\item \hyperref[stacks-cohomology-section-phantom]{Cohomologie des champs algébriques}
\item \hyperref[stacks-perfect-section-phantom]{Catégories dérivées de champs}
\item \hyperref[stacks-introduction-section-phantom]{Introduction aux champs algébriques}
\item \hyperref[stacks-more-morphisms-section-phantom]{Compléments sur les morphismes de champs}
\item \hyperref[stacks-geometry-section-phantom]{La géométrie des champs}
\end{enumerate}
Thèmes de théorie des espaces de modules
\begin{enumerate}
\setcounter{enumi}{107}
\item \hyperref[moduli-section-phantom]{Champs de modules}
\item \hyperref[moduli-curves-section-phantom]{Espaces de modules de courbes}
\end{enumerate}
Divers
\begin{enumerate}
\setcounter{enumi}{109}
\item \hyperref[examples-section-phantom]{Exemples}
\item \hyperref[exercises-section-phantom]{Exercices}
\item \hyperref[guide-section-phantom]{Guide bibliographique}
\item \hyperref[desirables-section-phantom]{Desiderata}
\item \hyperref[coding-section-phantom]{Style de programmation}
\item \hyperref[obsolete-section-phantom]{Obsolète}
\item \hyperref[fdl-section-phantom]{Licence de documentation libre GNU}
\item \hyperref[index-section-phantom]{Index généré automatiquement}
\end{enumerate}
\end{multicols}



\bibliography{my}
\bibliographystyle{amsalpha}

\end{document}
